\documentclass[runningheads]{llncs}
%
\usepackage[T1]{fontenc}
\usepackage{lineno}
\usepackage{graphicx}
\usepackage{amsmath}
\usepackage{amssymb}
\usepackage{amsthm}
\usepackage{xcolor}
\usepackage{textcomp}
\usepackage{listings}
\usepackage{courier}
\usepackage{algorithm}
\usepackage{algorithmic}
\usepackage{algpseudocode}
\usepackage[utf8]{inputenc}
\usepackage{geometry}
\geometry{left=2.5cm, right=2.5cm, top=2.5cm, bottom=2.5cm}

% Rust listings 颜色：防止 \color{rustkeyword} 未定义导致编译失败
\definecolor{rustkeyword}{RGB}{0,0,180}
\definecolor{rusttype}{RGB}{140,0,0}
\definecolor{rustcomment}{RGB}{0,110,0}
\definecolor{ruststring}{RGB}{160,0,160}

\newcommand{\N}{\mathbb{N}}
\newcommand{\Z}{\mathbb{Z}}

\newcommand{\Crate}{\mathcal{C}}        % crate
\newcommand{\SigmaC}{\Sigma(\Crate)}    % extracted env

\newcommand{\Cset}{\mathcal{C}}         % token color set (avoid clash with \Crate in prose)
\newcommand{\Qset}{\mathcal{Q}}
\newcommand{\Tset}{\mathcal{T}}

\newcommand{\Exp}{\text{Exp}}
\newcommand{\Type}{\text{Type}}
\newcommand{\GuardFun}{\text{Guard}}
\newcommand{\ActFun}{\text{Act}}
\newcommand{\Consume}{\text{Consume}}
\newcommand{\Produce}{\text{Produce}}

\newcommand{\Pop}{\mathsf{Pop}}
\newcommand{\Push}{\mathsf{Push}}

\newcommand{\dom}{\mathsf{dom}}
\newcommand{\ids}{\mathsf{Ids}}
\newcommand{\labs}{\mathsf{Labs}}
\newcommand{\Canon}{\mathsf{Canon}}
\newcommand{\Rename}{\mathsf{Rename}}
\newcommand{\TokObs}{\mathsf{Obs}}
\newcommand{\FreshOK}{\mathsf{FreshOK}}
\newcommand{\StackOK}{\mathsf{StackOK}}
\newcommand{\CapOK}{\mathsf{CapOK}}

\newcommand{\Join}{\mathbin{\uplus}}
\newcommand{\unify}{\mathrel{\dot{=}}}

\newcommand{\Cfg}{\mathsf{Cfg}}
\newcommand{\Fire}{\mathsf{Fire}}
\newcommand{\SolveGuard}{\mathsf{SolveGuard}}
\newcommand{\Closed}{\mathsf{Closed}}
\newcommand{\trace}{\mathsf{trace}}

\newcommand{\fvL}{\mathsf{fv}_L}
\newcommand{\restrict}{\upharpoonright}
\newcommand{\LamCut}[2]{#1\Downarrow #2}

\lstdefinelanguage{Rust}{
  keywords={fn, let, mut, struct, enum, impl, return, match, if, else, while, for, in, move, pub, use, mod, crate, ref, where, unsafe},
  keywordstyle=\color{rustkeyword}\bfseries,
  ndkeywords={String, Vec, Option, Result, u32, i32, bool, Self, Box, Rc, Arc},
  ndkeywordstyle=\color{rusttype}\bfseries,
  comment=[l]{//},
  commentstyle=\color{rustcomment}\itshape,
  stringstyle=\color{ruststring},
  sensitive=true,
  basicstyle=\ttfamily\small,
  breaklines=true,
  frame=single,
  numbers=left,
  numberstyle=\tiny\color{gray},
  captionpos=b,
  escapeinside={(*@}{@*)}, % 允许在代码中使用 LaTeX
}

\newtheorem{assumption}{Assumption}


\begin{document}
%
\title{Formally Verified Code Generation via Petri Nets for Rust’s Ownership and Lifetimes}
%
%\titlerunning{Abbreviated paper title}
% If the paper title is too long for the running head, you can set
% an abbreviated paper title here
%
\author{Kaiwen Zhang\inst{1}\orcidID{0000-1111-2222-3333} \and
Guanjun Liu\inst{1,2}\orcidID{1111-2222-3333-4444}}
%
\authorrunning{Kaiwen Zhang et al.}
% First names are abbreviated in the running head.
% If there are more than two authors, 'et al.' is used.
%
\institute{Tongji University, China 
\email{zhangkw@tongji.edu.cn}\\
\and
Jingangshan University, China\\
\email{liuguanjun@tongji.edu.cn}}
%
\maketitle              % typeset the header of the contribution
%
\begin{abstract}
% Safe Rust 强制执行编译时资源规则：所有权是仿射的，可变借用是排他性的，共享借用会冻结修改，引用受生命周期约束。在现代 crate 中，这些规则与泛型、trait 义务和关联类型相互作用，从而限制编译器接受哪些冗长且有状态的库调用序列。这里明确我们做的是签名级代码生成（不看函数体），并点出单测/模糊测试/修复三类场景。
Safe Rust enforces a compile-time resource discipline: ownership is affine, mutable borrows are exclusive, shared borrows freeze mutation, and references are constrained by lifetimes. In modern crates, these rules interact with generics, trait obligations, and associated types, thereby constraining which long, stateful
sequences of library calls are accepted by the compiler. We study signature-level compilation-admissible trace generation that given an extracted environment $\Sigma(\mathcal{C})$, synthesize self-contained client snippets that type-check and pass the borrow checker without inspecting function bodies, with applications to unit testing, fuzzing, and automated code repair.

% 本文利用Petri网，从语义角度出发，对这些编译约束进行了描述。我们定义了一个签名级的小步资源语义（SRS），它能够捕捉所有权转移、借用和再借用、生命周期约束以及特性/关联类型义务。基于$\Sigma(\mathcal{C})$，我们构建了一个下推着色Petri网（PCPN），其位置由允许的单态类型索引，并且其受保护的使能判断在触发时执行合一和义务检查。
This paper gives a semantics-first account of these compilation constraints using Petri nets. We define a signature-level small-step resource semantics (SRS) that captures ownership transfer, borrowing and reborrowing, lifetime outlives constraints, and trait/associated-type obligations. From $\Sigma(\mathcal{C})$ we construct a Pushdown Colored Petri Net (PCPN) whose places are indexed by admissible monomorphic types and whose guarded enabling judgement performs unification and obligation checking at firing time.

% 我们的主要成果是一个以互模拟形式表述的等价定理：由 Rust 签名层级资源规范所诱导的标记迁移系统，与 PCPN 的触发语义之间存在强互模拟关系。因此，一条单态化的 API 调用轨迹能被借用检查器接受的代码片段实现，当且仅当它对应于一条 PCPN 触发序列，且该序列最终到达一个closed configuration。在标准有限性假设下（即具体类型集合有限，且对令牌数量与借用深度设有显式上限），通过对状态空间进行canonical quotienting，可得到一个有限的可达图，从而支持完备的可达性分析、死锁检测、活性验证以及反例提取。
Our main result is a strong bisimulation (up to $\alpha$-renaming) between the SRS LTS and the PCPN firing LTS, implying that a monomorphized call trace is compilable iff it labels a firing sequence reaching a closed configuration.
To make exploration decidable, we introduce explicit bounds on token multiplicities and borrow depth, quotient by $\alpha$-equivalence via canonicalization, and obtain a finite reachability graph that serves as a search space for witness extraction and snippet emission. We further discuss strategy-guided exploration and instantiation policies that throttle endlessly applicable value producers and choose effective parameter/type bindings when multiple matches exist.
\keywords{Rust \and Colored Petri Nets \and Pushdown Nets \and Type Systems \and Bisimulation}
\end{abstract}

\section{Introduction}
% Rust 的安全子集通过所有权、借用和生命周期机制保障内存安全。除了防止内存错误之外，这套规范已成为 crate 作者编码 API 协议的一种实用手段：例如，子对象从父对象借用资源、会话（sessions）通过排他的&mut借用限制重入，而consuming methods则通过接收 self 强制实现一次性使用。此外，真实世界的 crate 大量依赖受 trait 约束和关联类型限制的泛型，因此合法的调用取决于在 crate 所处环境中能否满足签名中的约束条件。
Rust’s safe code enforces memory safety through ownership, borrowing, and lifetimes. Beyond preventing memory errors, this discipline has become a practical way for crate authors to encode API protocols that handles borrow from parents, sessions restrict re-entrancy through exclusive \texttt{\&mut} borrows, and consuming methods enforce one-shot usage by taking \texttt{self}. Moreover, real crates rely heavily on generics constrained by trait bounds and associated types, so admissible calls depend on solving signature constraints in the crate environment.

% 本文聚焦于signature-level的代码生成与分析：给定一个 crate C 及其提取出的签名/类型环境 Σ(C)，我们的目标是生成自包含的 Rust 代码片段，这些片段能以长而有状态的调用轨迹组合公开 API，且无需检查函数体。核心语义问题是轨迹可实现性:在 Σ(C) 环境下，哪些 API 调用轨迹可被一个可通过编译的 Rust 客户端实现？在 Rust 中，作用域不当的 move 操作、&mut下的非法别名，或未满足的 trait / 关联类型约束都会导致编译失败。因此，轨迹的可实现性主要由类型系统和借用检查器决定，而非运行时行为。
This paper focuses on the generation and analysis of signature-level code when given a crate $\mathcal{C}$ and its extracted type environment $\Sigma(\mathcal{C})$, and our objective is to generate self-contained Rust snippets that compose public APIs in long, stateful traces, without inspecting function bodies.
The core semantic question is trace realizability: Which API-call traces are realizable by a compilable Rust client under $\Sigma(\mathcal{C})$? In Rust, misdirected moves, illegal aliasing under \texttt{\&mut}, or unsatisfied trait/associated-type obligations block compilation, so trace realizability is governed primarily by the type system and the borrow checker rather than runtime behavior.

% Petri 网为推理追踪可实现性提供了一种自然的理论形式：资源对应于令牌，操作对应于转换，可实现性对应于可达性。然而，传统的位置/转换 (P/T) Petri 网缺乏足够的表达能力来捕捉 Rust 类型系统的三个关键方面：丰富的参数类型、由借用引起的能力转换以及由生命周期释放控制的权限返回。我们的关键洞见是，在签名层面，Safe Rust 可以抽象为一个有限的类型化资源转换器集合——每个转换器对应于由其函数签名指定的 API 调用——以及一组用于建模所有权转移和借用规则的小型固定结构操作。这些转换器不仅作用于类型化资源的多重集，而且还作用于一个作用域良好的借用结构中，该结构跟踪借用关系和生命周期。这一观察直接促使我们将其编码为下推着色 Petri 网 (PCPN)，它通过类型感知着色和类似栈的状态来丰富基础模型，从而表示作用域权限。
Petri nets offer a natural formalism for reasoning about trace realizability: resources correspond to tokens, operations to transitions, and realizability to reachability. However, classical place/transition (P/T) nets lack the expressiveness to capture three essential aspects of Rust's type system: rich parametric types, capability transformations induced by borrowing, and permission return governed by lifetime discharge. Our key insight is that, at the signature level, Safe Rust can be abstracted as a finite collection of typed resource transformers—each corresponding to an API call specified by its function signature—together with a small, fixed set of structural operations that model ownership transfer and borrowing discipline. These transformers act not only on a multiset of typed resources but also within a well-scoped loan structure that tracks borrow relationships and lifetimes. This observation directly motivates our encoding as a Pushdown Colored Petri Net (PCPN), which enriches the base model with both type-aware coloring and stack-like state to represent scoped permissions.

% 一个实用的程序生成器不应依赖于手动构建的标准库模型或外部依赖项。为了避免这种情况，我们直接从提取的编译环境 $\Sigma(\mathcal{C})$ 构建网络。编译器认为客户端代码可访问的每个构造函数、方法、trait 实现和关联类型相等性都被统一编码为一个源自其签名的转换。唯一不源自签名的值来源是综合框架提供的基本输入；所有用户定义和库定义的复合值都必须通过触发 $\Sigma(\mathcal{C})$ 中存在的转换来生成。这种自引导方法确保网络构建忠实地反映了编译器关于编译单元的封闭世界假设，从而保持了综合与实际编译可行性之间的一致性。
A practical program generator should not depend on a manually curated model of the standard library or external dependencies. To avoid this, we construct the net directly from the extracted compilation environment $\Sigma(\mathcal{C})$. Every constructor, method, trait implementation, and associated-type equality that the compiler deems accessible to client code is uniformly encoded as a transition derived from its signature. The only sources of values not originating from signatures are the primitive inputs supplied by the synthesis harness; all user-defined and library-defined composite values must be produced by firing transitions present in $\Sigma(\mathcal{C})$. This self-bootstrapping approach ensures that net construction faithfully mirrors the compiler’s closed-world assumption about the compilation unit, thereby maintaining alignment between synthesis and actual compilation feasibility.


This paper emphasizes semantic adequacy and analyzability:
\begin{enumerate}
  \item[1.]
  % 我们为 Safe Rust 定义了一个签名级资源语义（SRS），它精确地捕获了所有权转移、借用和再借用、生命周期约束以及特性和关联类型的推理——所有这些都无需检查函数体。这种语义将程序行为完全从类型签名和可见性信息中抽象出来。
  We define a signature-level resource semantics (SRS) for Safe Rust that precisely captures ownership transfer, borrowing and reborrowing, outlives constraints, and trait- and associated-type reasoning—all without inspecting function bodies. This semantics abstracts program behavior purely from type signatures and visibility information.

  \item[2.] 
  % 我们引入了一种直接从编译器提取的环境 $\Sigma(\mathcal{C})$ 构建的 crate 索引下推着色 Petri 网 (PCPN)，并在 SRS 的标记转换系统和 PCPN 的触发语义之间建立了强双模拟关系。由此得出一个当且仅当的刻画：当且仅当一个跟踪对应于 PCPN 执行达到闭合配置时，该跟踪才能被 Rust 的借用检查器编译。
  We introduce a crate-indexed Pushdown Colored Petri Net (PCPN) constructed directly from the compiler-extracted environment $\Sigma(\mathcal{C})$, and establish a strong bisimulation between the labeled transition system of the SRS and the firing semantics of the PCPN. This yields an if-and-only-if characterization: a trace is compilable by Rust’s borrow checker exactly when it corresponds to a PCPN execution reaching a closed configuration.
  
  \item[3.]
  % 我们通过统一的判断来形式化转换启用，该判断同时解决了类型-生命周期统一性、特征义务满足性、能力可行性和借用堆栈纪律问题。此判断作为每个 PCPN 转换的语义守卫，确保只允许借用检查器允许的步骤。
  We formalize transition enabling via a unified judgment that simultaneously resolves type–lifetime unification, trait obligation satisfaction, capability feasibility, and loan-stack discipline. This judgment serves as the semantic guard for each PCPN transition, ensuring that only borrow-checker–admissible steps are permitted.

  \item[4.] 
  % 在对标记重数和借用深度施加明确限制的条件下，我们通过规范化对状态进行模等价性运算，构建了一个有限可达性图 $\widehat{G}_{B,D}$。我们将 $\widehat{G}_{B,D}$ 视为一个用于代码生成的结构化搜索空间：见证触发序列会被回溯并生成语法有效的 Rust 代码片段。为了使搜索更具实用性，我们提出了一种策略引导的启发式方法，该方法 (i) 抑制对类常量值生成器的无限制应用，以及 (ii) 当多个候选者满足启用条件时，优先考虑高收益的参数和类型实例化。
  under explicit bounds on token multiplicities and borrow depth, we construct a finite reachability graph $\widehat{G}_{B,D}$ by quotienting states modulo equivalence through canonicalization. We treat $\widehat{G}_{B,D}$ as a structured search space for code generation: witness firing sequences are backtracked and emitted as syntactically valid Rust snippets. To make search practical, we propose strategy-guided heuristics that suppress unbounded application of constant-like value producers and prioritize high-yield parameter and type instantiations when multiple candidates satisfy the enabling judgment.
\end{enumerate}

Section~\ref{sec:background} reviews the Rust features and Petri-net preliminaries needed for the formal development. 
Section~\ref{sec:method} defines the PCPN construction, guard judgements, and the bisimulation theorem.
Section~\ref{sec:analysis} develops bounded reachability and classical analyses on the finite quotient.
An application layer that emits witnesses as Rust snippets is orthogonal and can be treated as a derived artifact.

\section{Background}
\label{sec:background}

\subsection{Rust Ownership, Borrowing, and Lifetimes}
\label{sec:background:rust}
% Rust 通过一套基于所有权与借用的静态机制来保障内存安全。任意时刻，每个非 Copy 类型的值都恰好有一个所有者；将该值赋给另一个变量通常会触发所有权的转移（即“移动”）。当所有者离开其作用域时，该值会被确定性地析构（drop），其所占用的资源随之释放。借用机制允许在不转移所有权的前提下临时访问某个值。Rust 区分两种引用形式：共享引用（&T）提供只读访问权限，并允许多个别名共存；可变引用（&mut T）则允许修改所指向的值，但必须保证独占性——即在其存活期间，不能存在任何其他共享或可变引用。这一限制通常被概括为“别名互斥可变性”（aliasing XOR mutability）原则：在程序执行的任一时刻，要么存在多个活跃的共享借用，要么存在唯一一个可变借用，二者不可同时成立。
Rust enforces memory safety through a static discipline grounded in ownership and borrowing. Every value of a non-Copy type has a unique owner at any given time; assignment to another variable typically transfers ownership via a move. When the owner’s scope ends, the value is deterministically dropped, and its associated resources are reclaimed. Borrowing provides temporary access to a value without transferring ownership. Rust distinguishes two forms of references: shared references $(\&T)$, which grant read-only access and may be freely aliased, and mutable references $(\&mut\ T)$, which permit mutation but must be uniquely held—no other references, shared or mutable, may coexist with a mutable borrow. This constraint is commonly encapsulated by the principle of aliasing XOR mutability: at any program point, either multiple shared borrows are active, or exactly one mutable borrow exists, but never both.

% 生命周期是编译期的概念，用于界定引用的有效范围。Rust 要求引用的生命周期不得超出其所指向值的存活时间；一旦借用结束，相应的访问权限即归还给所有者。在现代 Rust 中，借用的实际范围通常通过非词法生命周期（Non-Lexical Lifetimes, NLL）机制进行推断，这意味着借用的终止点由数据流分析决定，而非仅依赖于语法作用域。尽管这使得借用边界在源码表层不再显式可见，但它们对于判定后续操作（尤其是可变操作）是否合法仍至关重要。
Lifetimes, as compile-time annotations, delimit the regions over which references remain valid. A reference is required not to outlive the value it points to, and borrowing permissions are relinquished to the owner once the borrow’s lifetime ends. In contemporary Rust, the extents of borrows are often inferred using non-lexical lifetimes (NLL), meaning that the precise termination point of a borrow is determined by data-flow analysis rather than syntactic scope alone. Although this makes borrow boundaries less apparent from surface syntax, they remain critical for verifying the legality of subsequent operations, particularly mutations.

% 在签名级程序生成任务中，这些规则对合法调用序列施加了严格的语义约束。生成器必须精确维护每个值的状态信息——包括其是否仍被持有、是否已被移动、当前是否处于共享借用状态，或是否正被独占的可变借用所占用——并确保所生成的所有引用均遵守其来源约束，即不得在其所源自的值生命周期结束后继续使用。
For signature-level program generation, these rules impose nontrivial semantic constraints on valid call sequences. A generator must maintain an accurate model of each value’s state—whether it is still owned, has been moved, is currently subject to shared borrowing, or is exclusively mutably borrowed—and must ensure that all generated references respect their provenance, i.e., do not outlive the values they originate from.

% Rust 的 API 通常具有泛型性，函数常对类型参数进行量化，并通过 trait 约束加以限制，这些约束可直接出现在参数声明中，也可通过 where 子句表达。例如，函数签名 fn f<T: Trait>(x: T) -> U 仅允许将 T 实例化为实现了 Trait trait 的具体类型。从自动化程序生成的角度来看，调用此类多态函数需满足两个阶段的要求：首先，为泛型参数选择满足所有 trait 约束的具体类型；其次，构造出类型与实例化后签名相匹配的实参值。关键在于，若第一阶段失败——即不存在满足约束的有效类型实例化——则所生成的调用将无法通过类型检查，无论其运行时行为如何均无济于事。
Rust APIs are frequently generic, with functions commonly quantifying over type parameters subject to trait-based constraints, which may be expressed directly in parameter bounds or via where clauses. For instance, a function signature such as \texttt{fn f<T: Trait>(x: T) -> U} admits only those instantiations of T for which the type implements the trait Trait. From the standpoint of automated program generation, invoking such a polymorphic function entails a two-phase obligation: first, selecting concrete types that satisfy all declared trait bounds for the generic parameters; second, synthesizing argument values whose types conform to the resulting monomorphized signature. Crucially, if the first phase fails such that no valid instantiation exists that meets the required constraints, the generated call cannot pass type checking, irrespective of any consideration of dynamic semantics or runtime behavior.

\subsection{Signature-Induced Resource Semantics}
\label{sec:background:srs}
% 我们提出一种签名诱导的资源语义（Signature-Induced Resource Semantics, SRS），它精确刻画了代码生成所依赖的约束条件——这些约束完全从函数签名中推导得出，包括类型正确性、所有权转移、借用互斥性、生命周期释放，以及泛型义务的满足。值得注意的是，SRS 不解释函数体；每个函数调用被抽象为一个由签名上下文 Σ(C) 所允许的原子状态转换器。
We introduce a signature-induced resource semantics (SRS) that precisely captures the set of constraints essential for code generation as inferred solely from function signatures: namely, type correctness, ownership transfer, borrow exclusivity, lifetime discharge, and satisfaction of generic obligations. Notably, SRS abstracts over function bodies; each call is modeled as an atomic state transformer admitted by the signature context $\Sigma(\mathcal{C})$.

% 该语义中的一个配置是一个二元组 ⟨M,S⟩，其中 M 是在能力索引位置上的着色标记，S 是一个借用栈。此借用栈的结构与后文引入的参数化着色 Petri 网模型中的栈一致。M 表示当前可用的带类型资源及其关联的能力，例如可变性或借用权限。栈 S 则记录所有尚未释放的活跃借用，这些借用必须在其作用域结束前被正确归还。我们使用xx 表示一个由标签标识的合法单步转换。标签涵盖结构性动作，包括创建借用、结束借用、析构值和字段投影，也包括外部 API 调用。一个转换是否合法通过判断式规则定义，其依据与编译时验证机制一致。该机制包含类型与生命周期的合一、$\langle M,S\rangle$ 中 trait 与生命周期义务的验证，以及能力纪律与借用栈纪律的强制执行。因此，该语义在设计上是保守的，并与 Rust 编译器在签名层面的检查保持对齐
A configuration in this semantics is a pair $\langle M,S\rangle$ where $M$ is a colored marking over capability-indexed places and $S$ is a loan stack. The structure of $S$ matches that of the stack used in the PCPN introduced later. Intuitively $M$ represents the set of currently available typed resources together with their associated capabilities such as mutability or borrowing permissions. The stack $S$ records all outstanding borrows that must be discharged before their scopes end. We write
\begin{equation*}
    \Sigma(\mathcal{C}) \vdash \langle M,S\rangle \xRightarrow{l} \langle M',S'\rangle
\end{equation*}
to denote a single admissible transition labeled by $l$. Labels range over structural actions including borrow end-borrow drop and projection as well as external API calls. Admissibility of a transition is defined judgementally by the same constraint-solving discipline that underlies compile-time validation. This discipline encompasses type–lifetime unification verification of trait and lifetime obligations in $\Sigma(\mathcal{C})$ and enforcement of both capability discipline and loan stack discipline. As a result the semantics is deliberately conservative and aligned with the Rust compiler’s signature-level checks.

% 随后将证明，该语义所诱导的转换系统与参数化着色 Petri 网的触发语义在资源的 α-重命名下互模拟。这一等价性保证了两种形式体系均可作为生成程序中执行轨迹可实现性的可靠语义基础。
We later establish that the transition system induced by this semantics is bisimilar to the firing semantics of the PCPN modulo $\alpha$-renaming of resources. This equivalence guarantees that either formalism can serve as a sound semantic foundation for reasoning about trace realizability in generated programs.

% 为了解决从网到 Rust 语法的鸿沟，我们固定一个极小的语法片段，确保生成器输出的代码始终属于该片段；由此，语法正确性成为构造性不变量。
\subsection{A Borrow-Scoped Normal Form (BSNF) Core Rust Fragment}
\label{sec:background:core}

To address the Rust syntax gap, we fix a minimal syntactic fragment in which our emitter always produces code. This ensures syntactic correctness as a trivial invariant: emitted programs are by construction well-formed instances of the grammar.

% BSNF 的语句语法仅包含显式借用、投影、函数调用与显式释放操作，排除控制流与复杂表达式。
We adopt a restricted statement grammar:
\begin{align*}
\textsf{stmt}~::=&~
\textsf{let }v\textsf{ = } \textsf{Call}(f,\bar v);\ \mid\
\textsf{let }r\textsf{ = \&}v;\ \mid\
\textsf{let }r\textsf{ = \&mut }v;\ \mid\\
&\textsf{let }v^{\prime}\textsf{ = DerefCopy}(r);\ \mid\
\textsf{let }r^{\prime}\textsf{ = ReborrowShr}(r);\ \mid\
\textsf{let }r^{\prime}\textsf{ = ReborrowMut}(r);\ \mid\\
&\textsf{let }v^{\prime}\textsf{ = ProjMove}(v,f);\ \mid\
\textsf{let }r^{\prime}\textsf{ = ProjRef}(r,f);\ \mid\\
&\textsf{drop}(v);\ \mid\ \textsf{drop}(r);\\
\textsf{prog}~::=&~ \textsf{stmt}_1;\ \cdots\ ;\textsf{stmt}_n.
\end{align*}

% 函数调用是单态化的，可以以统一函数调用语法 (UFCS) 或方法形式出现。借用会将结果引用绑定到一个新的变量。显式的drop语句确定性地终止借用，避免依赖非词法生命周期启发式。投影仅限于基于移动的字段提取和通过引用重新借用，这与我们语义中建模的操作完全匹配。如果每个变量出现之前都有其唯一的绑定，并且所有变量名都不同，则 BSNF 程序在语法上是良好的。为了与我们的下推编码相匹配，我们限制 BSNF 见证必须被很好地括起来，即引用变量的显式drop必须遵循后进先出的顺序（相对于借用创建而言）。相对于完整的 Rust NLL，此限制较为保守，但足以满足见证的发出，并保持借用纪律的下推，从而在显式边界下进行商运算后能够进行有限状态分析。
Function calls are monomorphized and may appear in Uniform Function Call Syntax (UFCS) or method form. Borrows bind the resulting reference to a fresh variable.  
Explicit \textsf{drop} statements terminate borrows deterministically, avoiding reliance on non-lexical lifetime heuristics. Projections are limited to move-based field extraction and reborrowing through references, precisely matching the operations modeled in our semantics. A BSNF program is syntactically well-formed if every variable occurrence is preceded by its unique binding and all variable names are distinct. To match our pushdown encoding, we restrict BSNF witnesses to be well-bracketed that the explicit \textsf{drop} of reference variables must follow a last-in-first-out order with respect to borrow creation. This restriction is conservative relative to full Rust NLL, but it suffices for witness emission and keeps the loan discipline pushdown, enabling finite-state analysis after quotienting under explicit bounds.

% 编译器接受性被定义为程序在 SRS 下逐条语句可执行，即满足类型、义务、能力与栈约束。本工作的目标并非重形式化 rustc，而是提供一个以语义为中心、由签名驱动的编译约束刻画，足以支持见证生成。
We define compiler acceptance for BSNF programs via the judgement
\begin{equation*}
    \Sigma(\mathcal{C}) \vdash_{\textsf{BSNF}} \textsf{prog}:\textsf{ok},
\end{equation*}
which holds when each statement in sequence is admissible under the signature-induced resource semantics, that is when unification succeeds, all trait and outlives obligations are satisfied, capability constraints are met, and the loan stack evolves consistently. This formulation isolates the aspects we model signature-level typing and borrow checking from those we deliberately abstract away, such as function bodies and runtime behavior. The goal is not to re-formalize the entirety of rustc, but to provide a semantics-first, signature-driven characterization of compilation constraints relevant to code generation.  
Our BSNF fragment is sufficient for witness emission, and the SRS is conservative and aligned with the Rust compiler by construction.


\section{Pushdown Colored Petri Nets for Rust Signature-Level Generation}
\label{sec:method}

\subsection{Problem Setting}
\label{sec:method:setting}
% 设 $\mathcal{C}$ 为一个 Rust crate。我们可以从编译器前端提取出一个有限的可调用项集合 $\mathcal{F}$，其中包含自由函数与方法，并附带其完整签名。每个签名明确指定输入与输出类型、引用的可变性、显式的生命周期参数及其 outlives 约束（若存在），以及带有 trait 约束和关联类型约束的泛型参数。
Let $\mathcal{C}$ be a Rust crate. From the compiler front-end we can extract a finite set $\mathcal{F}$ of callable items comprising free functions and methods together with their full signatures. Each signature specifies input and output types the mutability of references any explicit lifetime parameters along with their outlives constraints and generic parameters equipped with trait bounds and associated-type constraints.

% 我们的目标是生成能够通过 Rust 编译器类型检查的代码片段，这些片段通过调用 $\mathcal{F}$ 中的项来构造类型正确的值并执行合法的调用序列。我们有意忽略函数体的内部实现细节；相反，每个调用被建模为一个原子转换，其编译时合法性完全由签名层面所体现的所有权、借用、生命周期及泛型义务等约束所决定。
Our objective is to generate Rust code snippets that are guaranteed to compile and that construct well-typed values by executing valid sequences of calls to items in $\mathcal{F}$. We deliberately abstract away from the internal implementation of function bodies; instead every call is modeled as an atomic transformation whose compile-time validity is determined entirely by the ownership borrowing lifetime and generic obligations as manifested at the signature level.

\subsection{Petri Net Preliminaries}
\label{sec:method:pn}
\begin{definition}[Net and Petri Net]
\label{def:pn}
A net is a $3$-tuple $N=(P,T,F)$, where $P$ is a finite set of places, $T$ is a finite set of transitions,
$F\subseteq (P\times T) \cup (T\times P)$ is a finite set of arcs, and $P\cap T=\emptyset$.
The pre-set and post-set of a node $x \in\{P \cup T\}$ are defined as
$^{\bullet}x=\{y\in P\cup T\mid (y,x)\in F\}$ and $x^{\bullet}=\{y\in P\cup T\mid (x,y)\in F\}$, respectively. 
\end{definition}
 A marking of $N=(P,T,F)$ is a mapping $M:P\rightarrow \mathbb{N}$ where $M(p)$ is the number of tokens in place $p$. A net $N$ with an initial marking $M_{0}$ is called a Petri Net and denoted as $(N,M_{0})$. Transition $t$ is enabled at $M$ if $\forall p\in {}^{\bullet}t: M(p)>0$, denoted as $M[t\rangle$. A marking $M$ is a deadlock marking if no transition is enabled at $M$. In this paper, it is assumed that the weight of each directed arc is $1$.

Firing an enabled transition $t$ yields a new marking $M^{\prime}$:
$M^{\prime}(p)=M(p)-1$ if $p\in {}^{\bullet}t\backslash t^{\bullet}$;
$M^{\prime}(p)=M^{\prime}(p)+1$ if $p\in t^{\bullet}\backslash {}^{\bullet}t$;
and otherwise $M^{\prime}(p)=M(p)$. This is denoted as $M[t\rangle M^{\prime}$. A marking $M_{k}$ is reachable from $M$ if $M_{k}=M$ or there exists a non-empty transition sequence $\sigma=t_{1}t_{2}\dots t_{k}$ such that $M[t_{1}\rangle M_{1}[t_{2}\rangle \dots M_{k-1}[t_{k}\rangle M_{k}$ (abbreviated as $M[\sigma \rangle M_{k}$). All markings reachable from $M_{0}$ in a Petri net is written as $\mathcal{M}$.


\subsection{Pushdown Colored Petri Nets (PCPN)}
\label{sec:method:pcpn}
% 设 $\mathsf{GTy}$ 为从 crate $\mathcal{C}$ 中提取出的单态 Rust 类型集合，$\mathsf{TVar}$ 为泛型类型变量集合，$\mathsf{LVar}$ 为生命周期变量集合，$\Sigma(\mathcal{C})$ 表示所提取的签名与类型环境。
Let $\mathsf{GTy}$ be the set of monomorphic Rust types extracted from $\mathcal{C}$. Let $\mathsf{TVar}$ be a set of generic type variables and $\mathsf{LVar}$ a set of lifetime variables. Let $\Sigma(\mathcal{C})$ denote the extracted signature and type environment.

% 我们采用如下类型表达式语言：基本类型来自 $\mathsf{GTy}$，类型变量来自 $\mathsf{TVar}$，支持类型构造、元组、切片、带生命周期和可变性的引用、关联类型投影、字段投影以及 trait 实现见证。
We use the following type-expression language:
\begin{align*}
\tau \in \mathsf{Ty} ::=~&
g \mid \alpha \mid \tau_0\langle \tau_1,\dots,\tau_n\rangle \mid (\tau_1,\dots,\tau_n) \mid [\tau] \mid \\
& \mathsf{Ref}^q_{\ell}(\tau) \mid \mathsf{Assoc}(\tau,\mathsf{Tr},A) \mid \mathsf{Field}(\tau,f) \mid \mathsf{Impl}(\tau,\mathsf{Tr}),
\end{align*}
where $g$ ranges over $\mathsf{GTy}$, $\alpha$ over $\mathsf{TVar}$, $q$ over $\{\mathsf{shr},\mathsf{mut}\}$,  
$\ell$ over $\mathsf{LVar}$, $\mathsf{Tr}$ over traits, $A$ over associated type names, and $f$ over field or projection paths.

% 我们引入一个保留的生命周期变量 $\ell_{\bullet}$，它不出现在任何提取出的签名中；所有内部生成的引用均使用 $\ell_{\bullet}$ 作为其类型中的生命周期，而实际区域标签通过估值函数 $\lambda(\ell_{\bullet})$ 携带。capability集合 $\mathsf{Cap}$ 包含四种状态：拥有、冻结、阻塞与值；基于此我们定义位置集合 $P$，每个位置由能力与规范化类型索引。
Let $\ell_{\bullet}\in \mathsf{LVar}$ be a reserved lifetime variable that does not occur in any extracted signature. All internally created references are typed with $\ell_{\bullet}$ and carry their actual region label via $\lambda(\ell_{\bullet})$. Let $\mathsf{Cap}=\{\mathsf{own},\mathsf{frz},\mathsf{blk}\}$.
We define the place set:
\begin{equation*}
   P ~=~ \{\, p_{\kappa,\hat\tau} \mid \kappa\in\mathsf{Cap},~ \hat\tau\in\widehat{\mathsf{Ty}} \,\}.
\end{equation*}

% 直观上，$p_{\mathsf{own},\hat\tau}$ 存储拥有权的值，$p_{\mathsf{frz},\hat\tau}$ 存储因共享借用而冻结的拥有值，$p_{\mathsf{blk},\hat\tau}$ 存储因可变借用而阻塞的拥有值，$p_{\mathsf{own},\hat\tau}$ 存储非拥有值，包括引用与见证。
Intuitively $p_{\mathsf{own},\hat\tau}$ stores all owned Rust values. The auxiliary places $p_{\mathsf{frz},\hat\tau}$ and $p_{\mathsf{blk},\hat\tau}$ store the same owned value temporarily marked as frozen/blocked due to outstanding borrows of its referent. In Rust, every \texttt{let}-binding owns its value even when that value is a reference, so we do not introduce a separatenon-owning value capability; borrowing changes permissions, not whether the value is owned.

% 设 $\mathbb{V}$ 为可数无限的值标识符集合（即新鲜的 Rust 变量名），$\mathbb{L}$ 为可数无限的区域标签集合。设 $\Theta$ 为有限类型代换集合，将类型变量映射到 $\mathsf{GTy}$；$\Lambda$ 为有限生命周期估值集合，将生命周期变量映射到 $\mathbb{L}$。
Let $\mathbb{V}$ be a countably infinite set of value identifiers corresponding to fresh Rust variable names, and $\mathbb{L}$ a countably infinite set of region labels. Let $\Theta$ be the set of finite type substitutions $\theta:\mathsf{TVar}\rightharpoonup \widehat{\mathsf{Ty}}$, and $\Lambda$ the set of finite lifetime valuations $\lambda:\mathsf{LVar}\rightharpoonup \mathbb{L}$.

We construct a PCPN
\begin{equation*}
     \mathcal{N}_{\mathcal{C}}=(P,T,W,\mathsf{Col},\Gamma,M_0,\mathsf{Guard},\mathsf{Out},\mathsf{Act}),
\end{equation*}
where $P$ is indexed by admissible ground types $\widehat{\mathsf{Ty}}$.

% 一个令牌颜色是一个三元组，包含值标识符、类型代换与生命周期估值。一个着色标记是从位置到有限彩色令牌多重集的映射。
A token color is
\begin{equation*}
     c ~=~ \langle v,\theta,\lambda\rangle \in \mathsf{Col} := \mathbb{V}\times\Theta\times\Lambda.
\end{equation*}
A colored marking is a mapping $M: P \to \mathbb{N}^{\mathsf{Col}}$, that is each place carries a finite multiset of colored tokens.

% 栈字母表 $\Gamma$ 包含三种操作：冻结、共享借用与可变借用，每种均携带值标识符与区域标签。
Let $\Gamma$ be a stack alphabet:
\begin{equation*}
    \Gamma ::= \mathsf{Freeze}(v)\mid \mathsf{Shr}(v_o,v_r,L)\mid \mathsf{Mut}(v_o,v_r,L),
\end{equation*}
where $L\in\mathbb{L}$ is a region label.

% 一个配置由一个着色标记与一个栈组成，其中栈顶位于左侧。弧的重数由一个具有有限支撑的权重函数建模，该函数定义在位置与变迁之间。对于变迁 $t$，我们用 $W(p,t)$ 与 $W(t,p)$ 分别表示其输入与输出重数。
A configuration is $\mathsf{Cfg}=\langle M,S\rangle$ where $S\in \Gamma^*$ with the top of the stack at the left. We model arc multiplicities by a weight function
\begin{equation*}
     W: (P\times T)\cup (T\times P)\to \mathbb{N}
\end{equation*}
with finite support, and let $F=\{(x,y)\mid W(x,y)>0\}$. For a transition $t$, we write $W(p,t)$ and $W(t,p)$ for its input and output multiplicities respectively.

% 在标记 $M$ 处对变迁 $t$ 的一个选择是一个多重集值函数，它从每个输入位置选取恰好所需数量的令牌，且这些令牌必须存在于当前标记中。
A choice for $t$ at marking $M$ is a multiset-valued function
\begin{equation*}
    \chi:P\to \mathbb{N}^{\mathsf{Col}}
    \quad\text{s.t.}\quad
    \forall p\in P:\ \chi(p)\subseteq M(p)\ \wedge\ |\chi(p)| = W(p,t). 
\end{equation*}

Let $\iota$ range over instantiations, containing at least:
\begin{equation*}
    \iota=(\theta^{*},\lambda^{*},\nu,\mu,\dots)
\end{equation*}
where $\theta^{*}\in\Theta$ and $\lambda^{*}\in\Lambda$ are the unified substitutions, and $\nu$ and $\mu$ provide any fresh value ids / fresh region labels needed by $\mathsf{Out}$. Other resolved information such as associated-type resolution may also be stored.

Each transition $t$ carries:
\begin{align*}
    & \mathsf{Guard}(t)(\langle M,S\rangle,\chi,\iota),\qquad \\
    & \mathsf{Out}(t)(\langle M,S\rangle,\chi,\iota):P\to \mathbb{N}^{(\mathsf{Col})},\qquad \\
    & \mathsf{Act}(t)\in\{\epsilon,\mathsf{push}(\bar\gamma),\mathsf{pop}(\bar\gamma)\}.
\end{align*}

\begin{definition} [Firing Rules]
\label{def:pcpn:firing}
A transition $t$ is enabled at $\mathsf{Cfg}$ if there exist a choice $\chi$ and instantiation record $\iota$ such that $\mathsf{Guard}(t)(\mathsf{Cfg},\chi,\iota)$ holds and any required pop-pattern matches the stack top. A firing rule is $\varphi=\langle t,\chi,\iota\rangle$.

Given $\varphi$, firing yields $\mathsf{Cfg}'=\langle M',S'\rangle$:
\begin{itemize}
\item Consume: $M^{-}(p)=M(p)\ominus \chi(p)$ for all $p$.
\item Produce: let $\pi=\mathsf{Out}(t)(\mathsf{Cfg},\chi,\iota)$, then $M'=M^{-}\oplus \pi$.
\item Stack: $S'$ is obtained by applying $\mathsf{Act}(t)$ to $S$.
\end{itemize}
We write $\mathsf{Cfg}[\varphi\rangle \mathsf{Cfg}'$. 
\end{definition}
A firing sequence is a finite sequence $\sigma=\varphi_1\cdots\varphi_n$ such that
$\mathsf{Cfg}_0[\sigma\rangle \mathsf{Cfg}_n$ where $\mathsf{Cfg}_0=\langle M_0,\epsilon\rangle$.
% 自举M0 只给出原始输入；复杂类型必须通过提供的构造函数/方法/trait 实现生成
$M_0$ seeds a finite multiset of primitive tokens that the synthesis harness may assume. Crucially, all composite values must be produced by firing transitions derived from $\Sigma(\mathcal{C})$. This avoids hand-modeling of the standard library: if $\Sigma(\mathcal{C})$ contains the relevant constructors/impls, the net can discover them; otherwise, the trace is unrealizable.

% ==========================
\subsection{Binding Elements and Guarded Enabling}
\label{sec:method:binding}
% 核心其实是将Rust类型系统中的类型约束处理机制，映射到PCPN的使能条件enabling condition中
% 在PCPN中，一个变迁被使能的条件是：存在一个binding element，即选择输入令牌并为过渡变量（transition variables）赋值，使守卫条件满足。在Rust类型系统的上下文中，这个binding同时完成两件事：泛型实例化：确定类型变量（如T）和生命周期变量（如'a）的具体值解决所有签名级obligation：如trait约束（T: Clone）、关联类型（Iterator::Item = i32）、生命周期约束（'a: 'b）。将Rust类型系统中三个步骤（类型统一、实例化完成、守卫求解）合并到一个框架中，避免在多个地方重复处理相同逻辑。
In PCPN, a transition is enabled by the existence of a binding element means a choice of input tokens together with an assignment to the transition variables that satisfies the guard. In our setting, this binding simultaneously determines generic instantiation such as type/lifetime variables and discharges all signature-level obligations, hence we present unification, instantiation completion, and guard solving in one place.

% Rust类型系统必须保证有限分支，像编译器不会陷入无限循环或爆炸式类型推导：
To keep enabling finite-branching, fix a finite ground-type pool $\mathsf{GTy}^{\dagger}\subseteq \mathsf{GTy}$. Let $\widehat{\mathsf{Ty}}$ be the finite closure of $\mathsf{GTy}^{\dagger}$ under the type constructors we model (e.g., tuples, slices, $\mathsf{Ref}^q_{\ell}(\cdot)$, and any finite crate-local projections). All place indices range over $\widehat{\mathsf{Ty}}$.

\subsubsection{Environment facts for traits and associated types}
\label{sec:method:assocfacts}

We assume the extracted environment $\Sigma(\mathcal{C})$ provides two finite, ground-level fact tables:
\begin{align*}
\mathsf{ImplFact}_{\Sigma} &\subseteq \widehat{\mathsf{Ty}}\times \mathsf{Trait},\\
\mathsf{AssocFact}_{\Sigma} &\subseteq \widehat{\mathsf{Ty}}\times \mathsf{Trait}\times \mathsf{AssocName}\times \widehat{\mathsf{Ty}}.
\end{align*}
Intuitively, $(\hat\tau,\mathsf{Tr})\in \mathsf{ImplFact}_{\Sigma}$ means the client can prove $\hat\tau:\mathsf{Tr}$;
and $(\hat\tau,\mathsf{Tr},A,\hat\tau_A)\in \mathsf{AssocFact}_{\Sigma}$ records the resolved associated type
$A$ of the unique impl of $\mathsf{Tr}$ for $\hat\tau$ when available to the client.

Rust coherence implies that for fixed $(\hat\tau,\mathsf{Tr},A)$ there is at most one $\hat\tau_A$ visible in the compilation unit. We capture this by the following property:
\begin{equation*}
    (\hat\tau,\mathsf{Tr},A,\hat\tau_1)\in \mathsf{AssocFact}_{\Sigma}\ \wedge\
(\hat\tau,\mathsf{Tr},A,\hat\tau_2)\in \mathsf{AssocFact}_{\Sigma}
\ \Rightarrow\ \hat\tau_1=\hat\tau_2.
\end{equation*}
Thus we may equivalently view associated-type resolution as a partial function
$\mathsf{AssocTy}_{\Sigma}(\hat\tau,\mathsf{Tr},A)=\hat\tau_A$.


\subsubsection{Obligations as guard atoms and entailment}
% 描述了Rust类型系统中如何处理trait约束、关联类型和生命周期的机制
Let $\mathsf{Obl}$ be the set of obligation atoms:
\begin{equation*}
    o \in \mathsf{Obl} ::=~ \tau:\mathsf{Tr}
    \ \mid\ 
    \mathsf{Assoc}(\tau,\mathsf{Tr},A)=\tau'
    \ \mid\
    (\ell_1:\ell_2).
\end{equation*}
An obligation set $\rho$ is a finite set of atoms. Trait bounds and associated-type equalities are treated as pure guard constraints.
% Rust将trait和关联类型推理完全作为guard处理，而不是在运行时生成额外的结构或转换
% 定义了义务原子（Obligation Atoms）来表示Rust类型系统中的约束，这些义务原子不是PCPN中的令牌，而是Rust类型系统中的约束，被映射为PCPN的守卫条件
Write $\sigma=(\theta,\lambda)$.  We define the entailment of instantiated obligations by the judgment
$\Sigma;S \vdash \sigma(o)$, with rules:
\[
\frac{
(\theta(\tau),\mathsf{Tr})\in \mathsf{ImplFact}_{\Sigma}
}{
\Sigma;S \vdash \sigma(\tau:\mathsf{Tr})
}
\ (\text{E-Trait})
\qquad
\frac{
(\theta(\tau),\mathsf{Tr})\in \mathsf{ImplFact}_{\Sigma}
\quad
\mathsf{AssocTy}_{\Sigma}(\theta(\tau),\mathsf{Tr},A)=\theta(\tau')
}{
\Sigma;S \vdash \sigma(\mathsf{Assoc}(\tau,\mathsf{Tr},A)=\tau')
}
\ (\text{E-Assoc})
\]

\[
\frac{
S \vdash \lambda(\ell_1)\sqsupseteq \lambda(\ell_2)
}{
\Sigma;S \vdash \sigma(\ell_1:\ell_2)
}
\ (\text{E-Life})
\]

and $\Sigma;S \vdash \sigma \models \rho$ holds iff every $o\in\rho$ is derivable.

\subsubsection{Unification as binding extraction with obligation emission}
% 统一（Unification）作为绑定提取，unification操作不仅返回实例化，还生成义务集合,TokObs(p,c)表示在库所p观察到的token，包含token 的类型和生命周期
Recall $\TokObs(p,c)=\langle \hat\tau,\LamCut{\lambda_c}{\hat\tau}\rangle$.
We refine unification to return both a forced instantiation $\sigma$ and a set of emitted obligations $\rho_u$:
\begin{equation*}
   \Sigma(\mathcal{C}) \vdash \tau \unify \TokObs(p,c) \Rightarrow (\sigma,\rho_u),
    \qquad where \sigma=(\theta,\lambda). 
\end{equation*}

We lift the join operation to pairs by
\begin{equation*}
  (\sigma_1,\rho_1)\Join(\sigma_2,\rho_2)\ \text{is defined iff}\ 
\sigma_1\parallel\sigma_2,\ \text{and then equals}\ (\sigma_1\Join\sigma_2,\ \rho_1\cup\rho_2).  
\end{equation*}
Joins over finite families are defined when all overlaps agree.

The following are the type unification rules:
\begin{equation*}
    \frac{}{~\Sigma \vdash g \unify \langle g,\lambda_t\rangle \Rightarrow ((\emptyset,\emptyset),\emptyset)~}
\ (\text{U-Ground}^{\rho})
\qquad
\frac{}{~\Sigma \vdash \alpha \unify \langle \hat\tau,\lambda_t\rangle \Rightarrow (([\alpha\mapsto \hat\tau],\emptyset),\emptyset)~}
\ (\text{U-TVar}^{\rho})
\end{equation*}

\begin{equation*}
    \frac{
\forall i.\ \Sigma \vdash \tau_i \unify \langle \hat\tau_i,\lambda_t\rangle \Rightarrow (\sigma_i,\rho_i)
\quad
(\sigma,\rho)= (\sigma_1,\rho_1)\Join\cdots\Join(\sigma_n,\rho_n)
}{
\Sigma \vdash \tau_0\langle \tau_1,\dots,\tau_n\rangle \unify
\langle \hat\tau_0\langle \hat\tau_1,\dots,\hat\tau_n\rangle,\lambda_t\rangle
\Rightarrow (\sigma,\rho)
}
\ (\text{U-STRUCT}^{\rho})
\end{equation*}

\begin{equation*}
  \frac{
\lambda_t(\ell')=L
\quad
\Sigma \vdash \tau \unify \langle \hat\tau,\lambda_t\rangle \Rightarrow (\sigma_0,\rho_0)
\quad
(\emptyset,[\ell\mapsto L]) \Join \sigma_0 = \sigma
}{
\Sigma \vdash \mathsf{Ref}^q_{\ell}(\tau) \unify \langle \mathsf{Ref}^q_{\ell'}(\hat\tau),\lambda_t\rangle
\Rightarrow (\sigma,\rho_0)
}
\ (\text{U-Ref}^{\rho})  
\end{equation*}

\begin{equation*}
    \frac{
\Sigma \vdash \tau \unify \langle \hat\tau,\lambda_t\rangle \Rightarrow (\sigma,\rho)
}{
\Sigma \vdash \mathsf{Field}(\tau,f) \unify \langle \mathsf{Field}(\hat\tau,f),\lambda_t\rangle \Rightarrow (\sigma,\rho)
}
\ (\text{U-Field}^{\rho})
\end{equation*}

Associated-type projections are not required to appear as place indices.
When a type scheme contains $\mathsf{Assoc}(\tau,\mathsf{Tr},A)$ and the observed token type is a ground type $\hat\tau$, we emit the obligations that force this projection to resolve to $\hat\tau$ under the chosen instantiation:
\begin{equation*}
    \frac{}{ \Sigma \vdash \mathsf{Assoc}(\tau,\mathsf{Tr},A) \unify \langle \hat\tau,\lambda_t\rangle \Rightarrow \Big((\emptyset,\emptyset),\ \{\tau:\mathsf{Tr},\ \mathsf{Assoc}(\tau,\mathsf{Tr},A)=\hat\tau\}\Big) }
\ (\text{U-AssocObs}^{\rho})
\end{equation*}


\subsubsection{Completion of partial bindings}
% 部分绑定的完成。unification操作返回的是强制的、可能导致部分的实例化。例如，unification可能只确定了部分类型变量，如T到i32，但未确定U到i64。
Unification yields a forced instantiation $\sigma_0=(\theta_0,\lambda_0)$.
We write $\sigma_0 \preceq \sigma^\star$ ($\sigma^\star$ extends $\sigma_0$) iff $\sigma_0$ and $\sigma^\star$ agree on $\dom(\sigma_0)$ and $\dom(\sigma_0)\subseteq \dom(\sigma^\star)$ componentwise.

Let $\bar\alpha(t)$ and $\bar\ell(t)$ be the type and lifetime variables quantified by $t$. A completed binding $\sigma^\star=(\theta^\star,\lambda^\star)$ for $t$ is any instantiation such that:
\begin{align*}
&\sigma_0 \preceq \sigma^\star,\\
&\dom(\theta^\star)\supseteq \bar\alpha(t)\ \wedge\ \forall \alpha\in \bar\alpha(t).\ \theta^\star(\alpha)\in \widehat{\mathsf{Ty}},\\
&\dom(\lambda^\star)\supseteq \bar\ell(t)\ \wedge\ \forall \ell\in \bar\ell(t).\ \lambda^\star(\ell)\in \mathbb{L}.
\end{align*}
% SolveGuard 的有限分支由 Ty 的有限性和有界的新标签预算保证 
% $\SolveGuard$ 的有限分支由 Ty 的有限性和有界的新标签预算保证。原则上，可以枚举所有补全 $\sigma^\star$ 以在选定的范围内恢复完备性。然而，在代码生成环境中，这可能会适得其反：许多类型可能满足相同的义务，探索所有类型可能会导致搜索过载。因此，我们允许采用一种类型选择策略，该策略对 Ty 中的候选基本类型进行排序：（例如，优先选择当前标记中已存在的类型；优先选择更小/更简单的类型；优先选择在 Crate 中具有更多可用构造函数/使用者的类型），并且仅探索每个调用点的前 k 个补全。
The finite branching of $\SolveGuard$ is ensured by the finiteness of $\widehat{\mathsf{Ty}}$ and bounded fresh-label budgets. In principle, one may enumerate all completions $\sigma^\star$ to recover completeness within the chosen bounds. In code-generation settings, however, this can be counterproductive: many types may satisfy the same obligations, and exploring all of them may overwhelm the search. We therefore allow a type-selection strategy that ranks candidate ground types in $\widehat{\mathsf{Ty}}$
(e.g., prefer types already present in the current marking; prefer maller/simpler types; prefer types with more available constructors/consumers in $\Sigma(\mathcal{C})$) and explores only the top-$k$ completions per call site.

\subsubsection{Guarded enabling judgement}
Let $\rho(t)$ be the obligation set extracted from the signature of transition $t$. Let $\Type_p(t,p)$ be the expected type scheme at input place $p$. We define guard solving as derivability of a binding judgement that merges obligations from the signature and those emitted by unification:
\begin{equation*}
    \SolveGuard(t,\langle M,S\rangle,\chi) ~:=~ \{\ \iota \mid \Sigma(\mathcal{C});\langle M,S\rangle \vdash_{\mathsf{solve}} t,\chi \Rightarrow \iota\ \}.
\end{equation*}

\begin{equation*}
\frac{
  \begin{aligned}
    &\chi \text{ is a choice for } t \text{ at } M \\
    &\forall p \in {}^{\bullet}t.\ \forall c \in \chi(p).\ 
      \Sigma \vdash \Type_p(t,p) \unify \TokObs(p,c) \Rightarrow (\sigma_{p,c},\rho_{p,c}) \\
    &(\sigma_0,\rho_u) = \Join_{p,c} (\sigma_{p,c},\rho_{p,c}) \\
    &\rho_{\mathrm{all}} = \rho(t)\ \cup\ \rho_u \\
    &\sigma^\star \text{ is a completed binding for } t \text{ extending } \sigma_0 \\
    &\Sigma;S \vdash \sigma^\star \models \rho_{\mathrm{all}} \\
    &\StackOK(t,S,\chi,\sigma^\star) \\
    &\CapOK(t,\langle M,S\rangle,\chi,\sigma^\star) \\
    &\FreshOK(\langle M,S\rangle,t,\chi,\sigma^\star,\nu,\mu)
  \end{aligned}
}{
  \Sigma;\langle M,S\rangle \vdash_{\mathsf{bind}} t,\chi \Rightarrow (\sigma^\star,\nu,\mu,\dots)
}
\ (\textsc{G-SOLVE})
\label{g-solve}
\end{equation*}
% 特征界限和相关类型等式仅通过 $\rho_{\mathrm{all}}$ 影响启用，因此它们不会引入额外的 Petri 网结构；它们是求解器侧的保护约束。
Trait bounds and associated-type equalities influence enabling only through $\rho_{\mathrm{all}}$, hence they do not introduce additional Petri-net structure; they are solver-side guard constraints.



\subsection{Constructing a Rust PCPN from $\mathcal{C}$}
\label{sec:method:construct}
We write $\mathsf{IsCopy}(\hat\tau)$ iff $\Sigma(\mathcal{C})\models \hat\tau:\mathsf{Copy}$.

We assume the extracted environment $\Sigma(\mathcal{C})$ induces a partial reduction function $\mathsf{Norm}_{\Sigma(\mathcal{C})}(\cdot)$ that eliminates well-formed associated-type projections $\mathsf{Assoc}(\hat\tau,\mathsf{Tr},A)$ into ground types using the unique trait impl information when available.
All places are indexed by normalized ground types in $\widehat{\mathsf{Ty}}$.

For each callable item $f\in\mathcal{F}$ with signature
\begin{equation*}
        \mathsf{sig}(f)=\langle \bar\alpha;\bar\ell;\ \tau_1,\dots,\tau_n \to \tau_o;\ \rho\rangle,
\end{equation*}
we add a transition $t_f$. Enabling of $t_f$ is defined by the guard judgement (\textsc{G-Solve}): it selects tokens, unifies types/lifetimes, and checks obligations $\rho$ including trait/associated-type constraints. Outputs are produced by $\mathsf{Out}$ using the normalized instantiated return type
$\mathsf{Norm}_{\Sigma(\mathcal{C})}((\theta,\lambda)(\tau_o))$.


\subsection{Ownership and Copy as Token Flow}
\label{sec:method:ownership}
Let $\pi_V(\langle v,\theta,\lambda\rangle)=v$, $\pi_\Theta(\cdot)=\theta$, and $\pi_\Lambda(\cdot)=\lambda$.
For a singleton multiset $\{c\}$, write $\mathsf{one}(\{c\})=c$.

\paragraph{(1) Move for non-Copy owned values.}
For each type $\tau$, add a transition $\mathsf{Move}_\tau$ as $T_{\tau}$ with:
\begin{equation*}
    W(p_{\mathsf{own},\tau},T_\tau)=1,
\end{equation*}
and no outputs. Define:
\begin{align*}
\mathsf{Guard}(T_\tau)(\langle M,S\rangle,\chi,\iota)
&~\equiv~ \neg \mathsf{IsCopy}(\tau),\\
\mathsf{Out}(T_\tau)(\langle M,S\rangle,\chi,\iota)
&~=~ \emptyset,\\
\mathsf{Act}(T_\tau)
&~=~ \epsilon.
\end{align*}

\paragraph{(2) By-value use for Copy owned values (self-loop).}
For each $\tau$, add $\mathsf{CopyUse}_\tau$ as $T_{\tau}$:
\begin{equation*}
    W(p_{\mathsf{own},\tau},T_\tau)=1,\qquad
    W(T_\tau,p_{\mathsf{own},\tau})=1.
\end{equation*}
Let $c=\mathsf{one}(\chi(p_{\mathsf{own},\tau}))$. Define:
\begin{align*}
\mathsf{Guard}(T_\tau)(\langle M,S\rangle,\chi,\iota)
&~\equiv~ \mathsf{IsCopy}(\tau),\\
\mathsf{Out}(T_\tau)(\langle M,S\rangle,\chi,\iota)
&~=~ [\,p_{\mathsf{own},\tau}\mapsto \{c\}\,],\\
\mathsf{Act}(T_\tau)
&~=~ \epsilon.
\end{align*}

\paragraph{(2$^{\star}$) Bounded duplication for Copy values.}
To preserve a one-to-one correspondence between live reference variables and stack frames in our well-bracketed fragment,
we do not provide duplication rules for reference-typed $\tau$.
Add $\mathsf{DupCopy}_\tau$ as $T_\tau$only when $\tau$ is Copy and not a reference type:
\begin{equation*}
    W(p_{\mathsf{own},\tau},T_\tau)=1,\qquad
    W(T_\tau,p_{\mathsf{own},\tau})=2.
\end{equation*}
Let $c=\langle v,\theta,\lambda\rangle=\mathsf{one}(\chi(p_{\mathsf{own},\tau}))$ and let $\nu(\mathsf{new})=v'$ be fresh.
Define:
\begin{align*}
\mathsf{Guard}(T_\tau)(\langle M,S\rangle,\chi,\iota)
&\equiv \mathsf{IsCopy}(\tau)\ \wedge\ \neg(\tau=\mathsf{Ref}^{q}_{\ell}(\cdot))\\
&\quad \wedge\ |M(p_{\mathsf{own},\tau})| < B_{\mathsf{copy}}(\tau)\ \wedge\ v'\notin \mathsf{Ids}(M,S),\\
\mathsf{Out}(T_\tau)(\langle M,S\rangle,\chi,\iota)
&= [\,p_{\mathsf{own},\tau}\mapsto \{\langle v,\theta,\lambda\rangle,\ \langle v',\theta,\lambda\rangle\}\,],\\
\mathsf{Act}(T_\tau)
&= \epsilon.
\end{align*}

\paragraph{(2$^{\prime}$) Duplication via explicit Clone.}
Unlike Copy, Clone creates a new value only via an explicit call. For each ground type $\tau$ such that
$\Sigma(\mathcal{C})\models \tau:\mathsf{Clone}$ and $\neg \mathsf{IsCopy}(\tau)$,
add $\mathsf{DupClone}_\tau$ as $T_{\tau}$:
\begin{equation*}
    W(p_{\mathsf{own},\tau},T_\tau)=1,\qquad
    W(T_\tau,p_{\mathsf{own},\tau})=2.
\end{equation*}
Let $c=\langle v,\theta,\lambda\rangle=\mathsf{one}(\chi(p_{\mathsf{own},\tau}))$
and let $\nu(\mathsf{new})=v'$ be fresh. Define:
\begin{align*}
\mathsf{Guard}(T_\tau)(\langle M,S\rangle,\chi,\iota)
&\equiv \Sigma(\mathcal{C})\models \tau:\mathsf{Clone}
\ \wedge\ \neg \mathsf{IsCopy}(\tau)\\
&\quad \wedge\ |M(p_{\mathsf{own},\tau})| < B_{\mathsf{clone}}(\tau)\ \wedge\ v'\notin \mathsf{Ids}(M,S),\\
\mathsf{Out}(T_\tau)(\langle M,S\rangle,\chi,\iota)
&= [\,p_{\mathsf{own},\tau}\mapsto \{\langle v,\theta,\lambda\rangle,\ \langle v',\theta,\lambda\rangle\}\,],\\
\mathsf{Act}(T_\tau)
&= \epsilon.
\end{align*}

\paragraph{(3) Drop only for non-borrowed non-reference values.}
We add drop transitions only for owned values that are not currently borrowed, and not reference-typed.
In Rust, dropping a reference value (\texttt{drop(r)}) is permitted but it semantically ends the borrow early;
we therefore model reference-dropping \emph{exclusively} by the end-borrow transitions below, and do not add
$\mathsf{DropOwn}$ for reference types.

\paragraph{(3$^{\prime}$) Drop owned non-reference values.}
For each ground type $\hat\tau$ such that $\hat\tau$ is not of the form $\mathsf{Ref}^{q}_{\ell}(\cdot)$,
add transition $\mathsf{DropOwn}_{\hat\tau}$ as $T_{\hat\tau}$:
\begin{equation*}
    W(p_{\mathsf{own},\hat\tau},T_{\hat\tau})=1,
\end{equation*}
with no outputs and $\mathsf{Act}=\epsilon$.


\paragraph{(4) First shared borrow from an owned value.}
Add transition $\mathsf{BorrowShrOwn}_\tau$ as $T_\tau$: \\
$W(p_{\mathsf{own},\tau},T_\tau)=W(T_\tau,p_{\mathsf{frz},\tau})=
W(T_\tau,p_{\mathsf{own},\mathsf{Ref}^{\mathsf{shr}}_{\ell_{\bullet}}(\tau)})=1.$

Let $c_o=\langle v_o,\theta,\lambda\rangle=\mathsf{one}(\chi(p_{\mathsf{own},\tau}))$ and let
$\nu(\mathsf{ref})=v_r$, $\mu(\mathsf{lab})=\ell$ be fresh.
Define:
\begin{align*}
\mathsf{Guard}(T_\tau)(\langle M,S\rangle,\chi,\iota)
&~\equiv~ v_r\notin \mathsf{Ids}(M,S)\ \wedge\ \ell\notin \mathsf{Labs}(M,S),\\
\mathsf{Out}(T_\tau)(\langle M,S\rangle,\chi,\iota)
&~=~ [\,p_{\mathsf{frz},\tau}\mapsto \{c_o\},\\
&\qquad\ \ p_{\mathsf{own},\mathsf{Ref}^{\mathsf{shr}}_{\ell_{\bullet}}(\tau)}\mapsto
\{\langle v_r,\theta,\lambda[\ell_{\bullet}\mapsto \ell]\rangle\}\,],\\
\mathsf{Act}(T_\tau)
&~=~ \mathsf{push}(\mathsf{Shr}(v_o,v_r,\ell)\cdot \mathsf{Freeze}(v_o)).
\end{align*}

\paragraph{(5) Additional shared borrow while already frozen.}
Add transition $\mathsf{BorrowShrFrz}_\tau$ as $T_\tau$: \\
$W(p_{\mathsf{frz},\tau},T_\tau)=W(T_\tau,p_{\mathsf{frz},\tau})=
W(T_\tau,p_{\mathsf{own},\mathsf{Ref}^{\mathsf{shr}}_{\ell_{\bullet}}(\tau)})=1.$
Let $c_o=\langle v_o,\theta,\lambda\rangle=\mathsf{one}(\chi(p_{\mathsf{frz},\tau}))$ and let
$\nu(\mathsf{ref})=v_r$, $\mu(\mathsf{lab})=\ell$ be fresh.
Define:
\begin{align*}
\mathsf{Guard}(T_\tau)(\langle M,S\rangle,\chi,\iota)
&~\equiv~ v_r\notin \mathsf{Ids}(M,S)\ \wedge\ \ell\notin \mathsf{Labs}(M,S),\\
\mathsf{Out}(T_\tau)(\langle M,S\rangle,\chi,\iota)
&~=~ [\,p_{\mathsf{frz},\tau}\mapsto \{c_o\},\\
&\qquad\ \ p_{\mathsf{own},\mathsf{Ref}^{\mathsf{shr}}_{\ell_{\bullet}}(\tau)}\mapsto
\{\langle v_r,\theta,\lambda[\ell_{\bullet}\mapsto \ell]\rangle\}\,],\\
\mathsf{Act}(T_\tau)
&~=~ \mathsf{push}(\mathsf{Shr}(v_o,v_r,\ell)).
\end{align*}

\paragraph{(6) Mutable borrow from an owned value.}
Add transition $\mathsf{BorrowMut}_\tau$ as $T_\tau$: \\
$W(p_{\mathsf{own},\tau},T_\tau)=W(T_\tau,p_{\mathsf{blk},\tau})=
W(T_\tau,p_{\mathsf{own},\mathsf{Ref}^{\mathsf{mut}}_{\ell_{\bullet}}(\tau)})=1.$
Let $c_o=\langle v_o,\theta,\lambda\rangle=\mathsf{one}(\chi(p_{\mathsf{own},\tau}))$ and let
$\nu(\mathsf{ref})=v_r$, $\mu(\mathsf{lab})=\ell$ be fresh.
Define:
\begin{align*}
\mathsf{Guard}(T_\tau)(\langle M,S\rangle,\chi,\iota)
&~\equiv~ v_r\notin \mathsf{Ids}(M,S)\ \wedge\ \ell\notin \mathsf{Labs}(M,S),\\
\mathsf{Out}(T_\tau)(\langle M,S\rangle,\chi,\iota)
&~=~ [\,p_{\mathsf{blk},\tau}\mapsto \{c_o\},\\
&\qquad\ \ p_{\mathsf{own},\mathsf{Ref}^{\mathsf{mut}}_{\ell_{\bullet}}(\tau)}\mapsto
\{\langle v_r,\theta,\lambda[\ell_{\bullet}\mapsto \ell]\rangle\}\,],\\
\mathsf{Act}(T_\tau)
&~=~ \mathsf{push}(\mathsf{Mut}(v_o,v_r,\ell)).
\end{align*}

\subsection{Lifetimes via Stack Binding}
\label{sec:method:lifetime}
\paragraph{(7) Ending a mutable borrow (restore owner).}
Add transition $\mathsf{EndMut}_\tau$ as $T_\tau$: \\
$W(p_{\mathsf{blk},\tau},T_\tau)=W(p_{\mathsf{own},\mathsf{Ref}^{\mathsf{mut}}_{\ell_{\bullet}}(\tau)},T_\tau)=
W(T_\tau,p_{\mathsf{own},\tau})=1.$

Let $c_o=\langle v_o,\theta_o,\lambda_o\rangle=\mathsf{one}(\chi(p_{\mathsf{blk},\tau}))$ and
$c_r=\langle v_r,\theta_r,\lambda_r\rangle=\mathsf{one}(\chi(p_{\mathsf{own},\mathsf{Ref}^{\mathsf{mut}}_{\ell_{\bullet}}(\tau)}))$.
Define:
\begin{align*}
\mathsf{Guard}(T_\tau)(\langle M,S\rangle,\chi,\iota)
& ~\equiv~ S = \mathsf{Mut}(v_o,v_r,\ell)\cdot S'
\ \wedge\ \theta_o=\theta_r
\ \wedge\ \LamCut{\lambda_o}{\tau} = \LamCut{\lambda_r}{\tau}. \\
\mathsf{Out}(T_\tau)(\langle M,S\rangle,\chi,\iota)
&~=~ [\,p_{\mathsf{own},\tau}\mapsto \{\langle v_o,\theta_o,\lambda_o\rangle\}\,],\\
\mathsf{Act}(T_\tau)
&~=~ \mathsf{pop}(\mathsf{Mut}(v_o,v_r,\ell)).
\end{align*}

\paragraph{(8) Ending a shared borrow (non-last).}
Add transition $\mathsf{EndShrKeep}_\tau$ as $T_\tau$: \\
$W(p_{\mathsf{frz},\tau},T_\tau)=W(p_{\mathsf{own},\mathsf{Ref}^{\mathsf{shr}}_{\ell_{\bullet}}(\tau)},T_\tau)=
W(T_\tau,p_{\mathsf{frz},\tau})=1.$

Let $c_o=\langle v_o,\theta,\lambda\rangle=\mathsf{one}(\chi(p_{\mathsf{frz},\tau}))$ and
$c_r=\langle v_r,\theta',\lambda'\rangle=\mathsf{one}(\chi(p_{\mathsf{own},\mathsf{Ref}^{\mathsf{shr}}_{\ell_{\bullet}}(\tau)}))$.
Define:
\begin{align*}
\mathsf{Guard}(T_\tau)(\langle M,S\rangle,\chi,\iota)
&~\equiv~ S=\mathsf{Shr}(v_o,v_r,\ell)\cdot S'
\ \wedge\ \theta=\theta'
\ \wedge\ \LamCut{\lambda}{\tau}=\LamCut{\lambda^{\prime}}{\tau} \\ 
&\quad \wedge \big(S'=\epsilon\ \vee\ \mathsf{top}(S')\neq \mathsf{Freeze}(v_o)\big),\\
\mathsf{Out}(T_\tau)(\langle M,S\rangle,\chi,\iota)
&~=~ [\,p_{\mathsf{frz},\tau}\mapsto \{\langle v_o,\theta,\lambda\rangle\}\,],\\
\mathsf{Act}(T_\tau)
&~=~ \mathsf{pop}(\mathsf{Shr}(v_o,v_r,\ell)).
\end{align*}

\paragraph{(9) Ending a shared borrow (last, restore owner).}
Add transition $\mathsf{EndShrLast}_\tau$ as $T_\tau$: \\
$W(p_{\mathsf{frz},\tau},T_\tau)=W(p_{\mathsf{own},\mathsf{Ref}^{\mathsf{shr}}_{\ell_{\bullet}}(\tau)},T_\tau)=
W(T_\tau,p_{\mathsf{own},\tau})=1.$

Let $c_o=\langle v_o,\theta,\lambda\rangle=\mathsf{one}(\chi(p_{\mathsf{frz},\tau}))$ and
$c_r=\langle v_r,\theta',\lambda'\rangle=\mathsf{one}(\chi(p_{\mathsf{own},\mathsf{Ref}^{\mathsf{shr}}_{\ell_{\bullet}}(\tau)}))$.
Define:
\begin{align*}
\mathsf{Guard}(T_\tau)(\langle M,S\rangle,\chi,\iota)
&~\equiv~ S=\mathsf{Shr}(v_o,v_r,\ell)\cdot \mathsf{Freeze}(v_o)\cdot S'
\ \wedge\ \theta=\theta'
\ \wedge\ \LamCut{\lambda}{\tau}=\LamCut{\lambda^{\prime}}{\tau},\\
\mathsf{Out}(T_\tau)(\langle M,S\rangle,\chi,\iota)
&~=~ [\,p_{\mathsf{own},\tau}\mapsto \{\langle v_o,\theta,\lambda\rangle\}\,],\\
\mathsf{Act}(T_\tau)
&~=~ \mathsf{pop}(\mathsf{Shr}(v_o,v_r,\ell)\cdot \mathsf{Freeze}(v_o)).
\end{align*}


\subsection{Composite Types and Field Re-Borrowing}
\label{sec:method:composite}

\paragraph{(10) Field projection by move.}
Let $\tau$ be a struct type with field $f:\tau_f$.
Add transition $\mathsf{ProjMove}_{\tau,f}$ as $T_{\tau,f}$: 
\begin{equation*}
W(p_{\mathsf{own},\tau},T_{\tau,f})=W(T_{\tau,f},p_{\mathsf{own},\tau_f})=1.
\end{equation*}
Let $c=\langle v,\theta,\lambda\rangle=\mathsf{one}(\chi(p_{\mathsf{own},\tau}))$ and let $\nu(\mathsf{fld})=v_f$ be fresh.
Define:
\begin{align*}
\mathsf{Guard}(T_{\tau,f})(\langle M,S\rangle,\chi,\iota)
&~\equiv~ v_f\notin \mathsf{Ids}(M,S),\\
\mathsf{Out}(T_{\tau,f})(\langle M,S\rangle,\chi,\iota)
&~=~ [\,p_{\mathsf{own},\tau_f}\mapsto \{\langle v_f,\theta,\lambda\rangle\}\,],\\
\mathsf{Act}(T_{\tau,f})
&~=~ \epsilon.
\end{align*}

\paragraph{(11) Field reborrow from a shared reference (Copy, keep parent).}
Add $\mathsf{ProjShr}_{\tau,f}$ as $T_{\tau,f}$:
\begin{equation*}
    W(p_{\mathsf{own},\mathsf{Ref}^{\mathsf{shr}}_{\ell}(\tau)},T_{\tau,f})=
    W(T_{\tau,f},p_{\mathsf{own},\mathsf{Ref}^{\mathsf{shr}}_{\ell}(\tau)})=
    W(T_{\tau,f},p_{\mathsf{own},\mathsf{Ref}^{\mathsf{shr}}_{\ell}(\tau_f)})=1.
\end{equation*}
Let $c_p=\langle v_p,\theta,\lambda\rangle=\mathsf{one}(\chi(p_{\mathsf{own},\mathsf{Ref}^{\mathsf{shr}}_{\ell}(\tau)}))$
and let $\nu(\mathsf{child})=v_c$ be fresh. Define:
\begin{align*}
\mathsf{Guard}(T_{\tau,f})(\langle M,S\rangle,\chi,\iota)
&~\equiv~ v_c\notin \mathsf{Ids}(M,S),\\
\mathsf{Out}(T_{\tau,f})(\langle M,S\rangle,\chi,\iota)
&~=~ [\,p_{\mathsf{own},\mathsf{Ref}^{\mathsf{shr}}_{\ell}(\tau)}\mapsto \{c_p\},\\
&\qquad\ \ p_{\mathsf{own},\mathsf{Ref}^{\mathsf{shr}}_{\ell}(\tau_f)}\mapsto \{\langle v_c,\theta,\lambda\rangle\}\,],\\
\mathsf{Act}(T_{\tau,f})
&~=~ \epsilon.
\end{align*}

\paragraph{(12) Field reborrow from a mutable reference (block parent ref explicitly).}
Add transition $\mathsf{ProjMut}_{\tau,f}$ as $T_{\tau,f}$: 
\begin{equation*}
    W(p_{\mathsf{own},\mathsf{Ref}^{\mathsf{mut}}_{\ell}(\tau)},T_{\tau,f})=
    W(T_{\tau,f},p_{\mathsf{blk},\mathsf{Ref}^{\mathsf{mut}}_{\ell}(\tau)})=
    W(T_{\tau,f},p_{\mathsf{own},\mathsf{Ref}^{\mathsf{mut}}_{\ell_{\bullet}}(\tau_f)})=1.
\end{equation*}
Let $c_p=\langle v_p,\theta,\lambda\rangle=\mathsf{one}(\chi(p_{\mathsf{own},\mathsf{Ref}^{\mathsf{mut}}_{\ell}(\tau)}))$
and let $\nu(\mathsf{child})=v_c$, $\mu(\mathsf{lab})=\ell_c$ be fresh. Define:
\begin{align*}
\mathsf{Guard}(T_{\tau,f})(\langle M,S\rangle,\chi,\iota)
&~\equiv~ v_c\notin \mathsf{Ids}(M,S)\ \wedge\ \ell_c\notin \mathsf{Labs}(M,S),\\
\mathsf{Out}(T_{\tau,f})(\langle M,S\rangle,\chi,\iota)
&~=~ [\,p_{\mathsf{blk},\mathsf{Ref}^{\mathsf{mut}}_{\ell}(\tau)}\mapsto \{c_p\},\\
&\qquad\ \ p_{\mathsf{own},\mathsf{Ref}^{\mathsf{mut}}_{\ell_{\bullet}}(\tau_f)}\mapsto
      \{\langle v_c,\theta,\lambda[\ell_{\bullet}\mapsto \ell_c]\rangle\}\,],\\
\mathsf{Act}(T_{\tau,f})
&~=~ \mathsf{push}(\mathsf{Mut}(v_p,v_c,\ell_c)).
\end{align*}

\paragraph{(13) Ending a mutable reborrow (restore parent ref from blk).}
Add transition $\mathsf{EndProjMut}_{\tau,f}$ as $T_{\tau,f}$:
\begin{equation*}
    W(p_{\mathsf{blk},\mathsf{Ref}^{\mathsf{mut}}_{\ell}(\tau)},T_{\tau,f})=
    W(p_{\mathsf{own},\mathsf{Ref}^{\mathsf{mut}}_{\ell_{\bullet}}(\tau_f)},T_{\tau,f})=
    W(T_{\tau,f},p_{\mathsf{own},\mathsf{Ref}^{\mathsf{mut}}_{\ell}(\tau)})=1.
\end{equation*}
Let $c_p=\langle v_p,\theta_p,\lambda_p\rangle=\mathsf{one}(\chi(p_{\mathsf{blk},\mathsf{Ref}^{\mathsf{mut}}_{\ell}(\tau)}))$ and
$c_c=\langle v_c,\theta_c,\lambda_c\rangle=\mathsf{one}(\chi(p_{\mathsf{own},\mathsf{Ref}^{\mathsf{mut}}_{\ell_{\bullet}}(\tau_f)}))$.
Define:
\begin{align*}
\mathsf{Guard}(T_{\tau,f})(\langle M,S\rangle,\chi,\iota)
&~\equiv~ S=\mathsf{Mut}(v_p,v_c,\ell_c)\cdot S',\\
\mathsf{Out}(T_{\tau,f})(\langle M,S\rangle,\chi,\iota)
&~=~ [\,p_{\mathsf{own},\mathsf{Ref}^{\mathsf{mut}}_{\ell}(\tau)}\mapsto \{\langle v_p,\theta_p,\lambda_p\rangle\}\,],\\
\mathsf{Act}(T_{\tau,f})
&~=~ \mathsf{pop}(\mathsf{Mut}(v_p,v_c,\ell_c)).
\end{align*}

\subsection{Dereference and Whole-Value Reborrowing}
\label{sec:method:deref}

\paragraph{(14) Copying out by dereference (requires Copy).}
For each ground type $\tau$, each qualifier $q\in\{\mathsf{shr},\mathsf{mut}\}$, and each lifetime variable $\ell$, add transition $\mathsf{DerefCopy}_{q,\tau,\ell}$ as $T_{q,\tau,\ell}$:
\begin{equation*}
  W(p_{\mathsf{own},\mathsf{Ref}^{q}_{\ell}(\tau)},T_{q,\tau,\ell})=1,\quad
  W(T_{q,\tau,\ell},p_{\mathsf{own},\mathsf{Ref}^{q}_{\ell}(\tau)})=1,\quad
  W(T_{q,\tau,\ell},p_{\mathsf{own},\tau})=1.
\end{equation*}
Let $c_r=\langle v_r,\theta,\lambda\rangle=\mathsf{one}(\chi(p_{\mathsf{own},\mathsf{Ref}^{q}_{\ell}(\tau)}))$
and let $\nu(\mathsf{new})=v$ be fresh. Define:
\begin{align*}
\mathsf{Guard}(T_{q,\tau,\ell})(\langle M,S\rangle,\chi,\iota)
&\equiv \mathsf{IsCopy}(\tau)\ \wedge\ v\notin \mathsf{Ids}(M,S),\\
\mathsf{Out}(T_{q,\tau,\ell})(\langle M,S\rangle,\chi,\iota)
&= [\,p_{\mathsf{own},\mathsf{Ref}^{q}_{\ell}(\tau)}\mapsto \{c_r\},\\
&\qquad p_{\mathsf{own},\tau}\mapsto \{\langle v,\theta,\lambda\rangle\}\,],\\
\mathsf{Act}(T_{q,\tau,\ell})
&= \epsilon.
\end{align*}

\paragraph{(15) Whole-value mutable reborrow from a mutable reference (block parent in blk).}
PCPNs model \texttt{let r2 = \&mut *r;} where \texttt{r: \&mut T}, keeping \texttt{r} usable after \texttt{r2} ends. For each $\tau$ and lifetime variable $\ell$, add transition $\mathsf{ReborrowMut}_{\tau,\ell}$ as $T_{\tau,\ell}$:
\begin{equation*}
  W(p_{\mathsf{own},\mathsf{Ref}^{\mathsf{mut}}_{\ell}(\tau)},T_{\tau,\ell})=1,\quad
  W(T_{\tau,\ell},p_{\mathsf{blk},\mathsf{Ref}^{\mathsf{mut}}_{\ell}(\tau)})=1,\quad
  W(T_{\tau,\ell},p_{\mathsf{own},\mathsf{Ref}^{\mathsf{mut}}_{\ell_{\bullet}}(\tau)})=1.
\end{equation*}
Let $c_p=\langle v_p,\theta,\lambda\rangle=\mathsf{one}(\chi(p_{\mathsf{own},\mathsf{Ref}^{\mathsf{mut}}_{\ell}(\tau)}))$
and let $\nu(\mathsf{child})=v_c$, $\mu(\mathsf{lab})=\ell_c$ be fresh. Define:
\begin{align*}
\mathsf{Guard}(T_{\tau,\ell})(\langle M,S\rangle,\chi,\iota)
&\equiv v_c\notin \mathsf{Ids}(M,S)\ \wedge\ \ell_c\notin \mathsf{Labs}(M,S),\\
\mathsf{Out}(T_{\tau,\ell})(\langle M,S\rangle,\chi,\iota)
&= [\,p_{\mathsf{blk},\mathsf{Ref}^{\mathsf{mut}}_{\ell}(\tau)}\mapsto \{c_p\},\\
&\qquad p_{\mathsf{own},\mathsf{Ref}^{\mathsf{mut}}_{\ell_{\bullet}}(\tau)}\mapsto
      \{\langle v_c,\theta,\lambda[\ell_{\bullet}\mapsto \ell_c]\rangle\}\,],\\
\mathsf{Act}(T_{\tau,\ell})
&= \mathsf{push}(\mathsf{Mut}(v_p,v_c,\ell_c)).
\end{align*}

Add the corresponding ending transition $\mathsf{EndReborrowMut}_{\tau,\ell}$ as $T_{\tau,\ell}$:
\begin{equation*}
  W(p_{\mathsf{blk},\mathsf{Ref}^{\mathsf{mut}}_{\ell}(\tau)},T_{\tau,\ell})=1,\quad
  W(p_{\mathsf{own},\mathsf{Ref}^{\mathsf{mut}}_{\ell_{\bullet}}(\tau)},T_{\tau,\ell})=1,\quad
W(T_{\tau,\ell},p_{\mathsf{own},\mathsf{Ref}^{\mathsf{mut}}_{\ell}(\tau)})=1.
\end{equation*}
Let $c_p=\langle v_p,\theta_p,\lambda_p\rangle=\mathsf{one}(\chi(p_{\mathsf{blk},\mathsf{Ref}^{\mathsf{mut}}_{\ell}(\tau)}))$ and
$c_c=\langle v_c,\theta_c,\lambda_c\rangle=\mathsf{one}(\chi(p_{\mathsf{own},\mathsf{Ref}^{\mathsf{mut}}_{\ell_{\bullet}}(\tau)}))$.
Define:
\begin{align*}
\mathsf{Guard}(T_{\tau,\ell})(\langle M,S\rangle,\chi,\iota)
&\equiv S=\mathsf{Mut}(v_p,v_c,\ell_c)\cdot S',\\
\mathsf{Out}(T_{\tau,\ell})(\langle M,S\rangle,\chi,\iota)
&= [\,p_{\mathsf{own},\mathsf{Ref}^{\mathsf{mut}}_{\ell}(\tau)}\mapsto \{\langle v_p,\theta_p,\lambda_p\rangle\}\,],\\
\mathsf{Act}(T_{\tau,\ell})
&= \mathsf{pop}(\mathsf{Mut}(v_p,v_c,\ell_c)).
\end{align*}

\paragraph{(16) Shared reborrow from a mutable reference (block parent in blk).}
This models \texttt{let s = \&*r;} and the coercion \texttt{\&mut T} $\rightsquigarrow$ \texttt{\&T}.
For each $\tau$ and lifetime variable $\ell$, add transition $\mathsf{ReborrowShrFromMut}_{\tau,\ell}$ as $T_{\tau,\ell}$:
\begin{equation*}
  W(p_{\mathsf{own},\mathsf{Ref}^{\mathsf{mut}}_{\ell}(\tau)},T_{\tau,\ell})=1,\quad
  W(T_{\tau,\ell},p_{\mathsf{blk},\mathsf{Ref}^{\mathsf{mut}}_{\ell}(\tau)})=1,\quad
  W(T_{\tau,\ell},p_{\mathsf{own},\mathsf{Ref}^{\mathsf{shr}}_{\ell_{\bullet}}(\tau)})=1.
\end{equation*}
Let $c_p=\langle v_p,\theta,\lambda\rangle=\mathsf{one}(\chi(p_{\mathsf{own},\mathsf{Ref}^{\mathsf{mut}}_{\ell}(\tau)}))$
and let $\nu(\mathsf{child})=v_c$, $\mu(\mathsf{lab})=\ell_c$ be fresh. Define:
\begin{align*}
\mathsf{Guard}(T_{\tau,\ell})(\langle M,S\rangle,\chi,\iota)
&\equiv v_c\notin \mathsf{Ids}(M,S)\ \wedge\ \ell_c\notin \mathsf{Labs}(M,S),\\
\mathsf{Out}(T_{\tau,\ell})(\langle M,S\rangle,\chi,\iota)
&= [\,p_{\mathsf{blk},\mathsf{Ref}^{\mathsf{mut}}_{\ell}(\tau)}\mapsto \{c_p\},\\
&\qquad p_{\mathsf{own},\mathsf{Ref}^{\mathsf{shr}}_{\ell_{\bullet}}(\tau)}\mapsto
      \{\langle v_c,\theta,\lambda[\ell_{\bullet}\mapsto \ell_c]\rangle\}\,],\\
\mathsf{Act}(T_{\tau,\ell})
&= \mathsf{push}(\mathsf{Mut}(v_p,v_c,\ell_c)).
\end{align*}

Add the corresponding ending transition $\mathsf{EndReborrowShrFromMut}_{\tau,\ell}$ as $T_{\tau,\ell}$:
\begin{equation*}
  W(p_{\mathsf{blk},\mathsf{Ref}^{\mathsf{mut}}_{\ell}(\tau)},T_{\tau,\ell})=1,\quad
  W(p_{\mathsf{own},\mathsf{Ref}^{\mathsf{shr}}_{\ell_{\bullet}}(\tau)},T_{\tau,\ell})=1,\quad
  W(T_{\tau,\ell},p_{\mathsf{own},\mathsf{Ref}^{\mathsf{mut}}_{\ell}(\tau)})=1.
\end{equation*}
Let $c_p=\langle v_p,\theta_p,\lambda_p\rangle=\mathsf{one}(\chi(p_{\mathsf{blk},\mathsf{Ref}^{\mathsf{mut}}_{\ell}(\tau)}))$ and
$c_c=\langle v_c,\theta_c,\lambda_c\rangle=\mathsf{one}(\chi(p_{\mathsf{own},\mathsf{Ref}^{\mathsf{shr}}_{\ell_{\bullet}}(\tau)}))$.
Define:
\begin{align*}
\mathsf{Guard}(T_{\tau,\ell})(\langle M,S\rangle,\chi,\iota)
&\equiv S=\mathsf{Mut}(v_p,v_c,\ell_c)\cdot S',\\
\mathsf{Out}(T_{\tau,\ell})(\langle M,S\rangle,\chi,\iota)
&= [\,p_{\mathsf{own},\mathsf{Ref}^{\mathsf{mut}}_{\ell}(\tau)}\mapsto \{\langle v_p,\theta_p,\lambda_p\rangle\}\,],\\
\mathsf{Act}(T_{\tau,\ell})
&= \mathsf{pop}(\mathsf{Mut}(v_p,v_c,\ell_c)).
\end{align*}

\subsection{Correctness: Step Correspondence and Bisimulation}
\label{sec:method:correctness}

We state correctness as an equivalence between SRS from Section~\ref{sec:background:srs}, and the PCPN firing semantics. 

Let $\mathsf{Lab}$ be the set of step labels, containing structural labels for ownership/borrowing/projection/drop actions and API-call labels $f@(\theta,\lambda)$ recording the called item and the resolved instantiation.
We write $\mathrm{lab}(\varphi)\in\mathsf{Lab}$ for the label of an instantiated firing $\varphi=\langle t,\chi,\iota\rangle$, defined by cases on $t$:
\begin{itemize}
  \item if $t=t_f$ is derived from an item $f\in\mathcal{F}$, then
  $\mathrm{lab}(\varphi)=f@(\theta,\lambda)$ where $(\theta,\lambda)$ is the unified instantiation stored in $\iota$;
  \item if $t$ is a structural schema such that borrow/end-borrow/projection/drop/move/clone, then
  $\mathrm{lab}(\varphi)$ is the corresponding structural label including its ground type parameters and projection path, if any.
\end{itemize}

As firings allocate fresh value ids and region labels, we quotient configurations by $\alpha$-renaming. A renaming $\rho=(\rho_V,\rho_L)$ is a pair of bijections on value identifiers and region labels. We write $\rho(\langle M,S\rangle)$ for its componentwise action, and define $\mathsf{Cfg}_1\equiv_\alpha \mathsf{Cfg}_2$ iff $\mathsf{Cfg}_2=\rho(\mathsf{Cfg}_1)$ for some $\rho$.

We relate SRS configurations and PCPN configurations by identity up to $\alpha$:
\begin{equation*}
    (\mathsf{Cfg}_{\mathrm{srs}},\mathsf{Cfg}_{\mathrm{pcpn}})\in\mathcal{R}
\quad\Longleftrightarrow\quad
\mathsf{Cfg}_{\mathrm{srs}}\equiv_\alpha \mathsf{Cfg}_{\mathrm{pcpn}}.
\end{equation*}

\begin{lemma}[Schema coverage and uniqueness]
\label{lem:schema-coverage}
For every SRS step label $l\in\mathsf{Lab}$, there exists a unique transition schema $t_l\in T$ such that any PCPN firing with label $l$ must have underlying transition $t_l$.
\end{lemma}
\noindent\textbf{Proof sketch.}
By construction of $\mathcal{N}_{\mathcal{C}}$: API labels $f@(\theta,\lambda)$ arise only from the transition $t_f$ added for $f\in\mathcal{F}$; structural labels arise only from their dedicated schema. Hence uniqueness holds by case analysis on $l$. \hfill $\square$

\begin{lemma}[Guard judgement characterizes enabling]
\label{lem:guard-characterizes}
For any transition $t$, configuration $\mathsf{Cfg}$, and choice $\chi$,
there exists an instantiation record $\iota$ such that $t$ is enabled at $\mathsf{Cfg}$ w.r.t.\ $\chi$ and $\iota$ iff $\iota\in\SolveGuard(t,\mathsf{Cfg},\chi)$ and the stack pop-pattern matches the stack top.
\end{lemma}
\noindent\textbf{Proof.}
Immediate from Definition~\ref{def:pcpn:firing} and the definition of $\SolveGuard$ in Section~\ref{g-solve}. \hfill $\square$

\begin{lemma}[$\alpha$-invariance of firing]
\label{lem:alpha-invariance-method}
If $\mathsf{Cfg}_1\equiv_\alpha \mathsf{Cfg}_2$, then for every enabled instantiated firing
$\varphi$ at $\mathsf{Cfg}_1$ there exists a corresponding enabled firing $\varphi'$ at $\mathsf{Cfg}_2$
such that:
\begin{equation*}
    \mathrm{lab}(\varphi)=\mathrm{lab}(\varphi')
    \quad\text{and}\quad
    \mathsf{Fire}(\mathsf{Cfg}_1,\varphi)\equiv_\alpha \mathsf{Fire}(\mathsf{Cfg}_2,\varphi').
\end{equation*}
And conversely.
\end{lemma}
\noindent\textbf{Proof sketch.}
Renaming preserves place indices, multiset multiplicities, and stack-frame shapes. All guard side conditions are stable under bijective renaming of ids/labels: equality/inequality tests are preserved, unification observes only types and restricted valuations, and freshness constraints can be transported by composing with $\rho$. Thus enabledness and the post-configuration are preserved up to $\alpha$. \hfill $\square$

\begin{lemma}[One-step correspondence]
\label{lem:one-step-corr}
For any configuration $\mathsf{Cfg}$ and label $l\in\mathsf{Lab}$,
\begin{equation*}
    \Sigma(\mathcal{C}) \vdash \mathsf{Cfg} \xRightarrow{l} \mathsf{Cfg}'
    \quad\Longleftrightarrow\quad
    \exists \varphi.\ \mathsf{Cfg}[\varphi\rangle \mathsf{Cfg}' \ \wedge\ \mathrm{lab}(\varphi)=l,
\end{equation*}
where $\xRightarrow{l}$ is the SRS step relation and $[\varphi\rangle$ is PCPN firing, both understood modulo $\alpha$-renaming of freshly allocated ids/labels.
\end{lemma}
\noindent\textbf{Proof sketch.}
($\Rightarrow$) By Lemma~\ref{lem:schema-coverage}, the label $l$ determines a unique schema $t_l$. SRS admissibility conditions for the step are exactly the constraints checked by guard solving: type/lifetime unification, obligation entailment, capability feasibility and stack discipline
(encoded in $\CapOK$ and $\StackOK$). Hence some $\iota\in\SolveGuard(t_l,\mathsf{Cfg},\chi)$ exists, and firing yields the same resource effect by construction of $\mathsf{Out}$ and $\mathsf{Act}$. Fresh-name differences are reconciled by $\alpha$-renaming (Lemma~\ref{lem:alpha-invariance-method}).

($\Leftarrow$) Conversely, any PCPN firing uses a schema $t_l$ whose guard judgement holds, thus the corresponding SRS rule is admissible and yields the same update on marking and stack, again up to $\alpha$. \hfill $\square$

\begin{theorem}[SRS--PCPN strong bisimulation up to $\alpha$]
\label{thm:bisim}
The relation $\mathcal{R}$ is a strong bisimulation up to $\alpha$ between the SRS LTS and the PCPN firing LTS.
\end{theorem}
\noindent\textbf{Proof.}
Assume $(\mathsf{Cfg}_{\mathrm{srs}},\mathsf{Cfg}_{\mathrm{pcpn}})\in\mathcal{R}$, i.e. they are $\alpha$-equivalent.
Let $\mathsf{Cfg}_{\mathrm{srs}} \xRightarrow{l} \mathsf{Cfg}'_{\mathrm{srs}}$ be any SRS step.
By Lemma~\ref{lem:one-step-corr}, there exists a PCPN firing $\varphi$ with label $l$ producing
$\mathsf{Cfg}'_{\mathrm{pcpn}}$ such that $\mathsf{Cfg}'_{\mathrm{srs}}\equiv_\alpha \mathsf{Cfg}'_{\mathrm{pcpn}}$,
hence $(\mathsf{Cfg}'_{\mathrm{srs}},\mathsf{Cfg}'_{\mathrm{pcpn}})\in\mathcal{R}$.
The reverse direction follows symmetrically using the same lemma.
Thus $\mathcal{R}$ is a strong bisimulation up to $\alpha$. \hfill $\square$

\begin{corollary}[Trace realizability iff closed reachability]
\label{cor:trace}
A monomorphized API-call trace is realizable by a compilable Rust client under $\Sigma(\mathcal{C})$
iff it is the projection of a PCPN firing sequence that reaches a closed configuration
(i.e., a configuration with empty loan stack).
\end{corollary}

\section{Bounded Reachability Analysis and Trusted Snippet Synthesis}
\label{sec:analysis}
This section focuses on constructing a finite reachability graph on top of the crate-indexed PCPN $\mathcal{N}_{\mathcal{C}}$ (Section~\ref{sec:method}) and using it as a search space for compilable witnesses. Because the raw state space is infinite such as unbounded fresh ids/labels, unbounded stack depth, and potentially unbounded monomorphization, all completeness claims in this section are relative to explicit bounds.


\subsection{Bounded firing semantics}
\label{sec:analysis:bounded}

Recall a configuration $\Cfg=\langle M,S\rangle$ (Definition~\ref{def:pcpn:firing}).
For a multiset $X\in\mathbb{N}^{\mathsf{Col}}$, write $|X|:=\sum_{c\in\mathsf{Col}}X(c)$.
Thus $|M(p)|$ is the total number of tokens in place $p$.

\paragraph{Bounds.}
Fix:
\begin{itemize}
  \item a finite ground-type universe $\widehat{\mathsf{Ty}}$ (already fixed by the net construction),
  \item a place budget $B:P\to\mathbb{N}$ bounding token multiplicities,
  \item a stack-depth bound $D\in\mathbb{N}$.
\end{itemize}
We write $\mathsf{BudOK}_{B,D}(\langle M,S\rangle)$ iff
$\forall p\in P.\ |M(p)|\le B(p)$ and $|S|\le D$.

\paragraph{Bounded step.}
Let $\varphi=\langle t,\chi,\iota\rangle$ be an instantiated firing rule.
We define a bounded step relation:
\[
\Cfg[\varphi\rangle_{B,D}\Cfg'
\quad\Longleftrightarrow\quad
\Cfg[\varphi\rangle\Cfg' \ \wedge\ \mathsf{BudOK}_{B,D}(\Cfg').
\]
Bounded reachability $\Cfg \Rightarrow^{*}_{B,D} \Cfg'$ is the reflexive-transitive closure of $\rangle_{B,D}$.

\subsection{$\alpha$-equivalence and canonicalization}
\label{sec:analysis:canon}

Fresh ids/labels introduced by firings yield infinitely many syntactically distinct but isomorphic configurations.
We therefore quotient by $\alpha$-renaming and explore only canonical representatives.

A renaming $\rho=(\rho_V,\rho_L)$ is a pair of bijections on value identifiers and region labels.
Write $\rho(\Cfg)$ for the componentwise action on markings and stacks.
Define $\Cfg_1\equiv_\alpha \Cfg_2$ iff $\Cfg_2=\rho(\Cfg_1)$ for some $\rho$.

\begin{definition}
\label{def:canon}
A canonicalization function $\Canon:\Cfg\to\Cfg$ returns a representative of each $\alpha$-equivalence class,
obtained by:
(i) renaming value ids and region labels to a fixed order induced by first occurrence (stack top-to-bottom,
then marking place-by-place), and
(ii) sorting token multisets within each place by a fixed total order on colors.
\end{definition}

\begin{lemma}[$\alpha$-invariance of bounded firing]
\label{lem:alpha-invariance-analysis}
If $\Cfg_1\equiv_\alpha \Cfg_2$ and $\Cfg_1[\varphi\rangle_{B,D}\Cfg_1'$,
then there exist $\varphi'$ and $\Cfg_2'$ such that
$\Cfg_2[\varphi'\rangle_{B,D}\Cfg_2'$, $\mathrm{lab}(\varphi)=\mathrm{lab}(\varphi')$,
and $\Cfg_1'\equiv_\alpha \Cfg_2'$.
\end{lemma}
\noindent\textbf{Proof sketch.}
Immediate from Lemma~\ref{lem:alpha-invariance-method}; the budget predicate is stable under renaming. \hfill $\square$

\begin{theorem}[Quotient soundness]
\label{thm:quotient-soundness}
For any configurations $\Cfg,\Cfg'$,
\[
\Cfg \Rightarrow^{*}_{B,D} \Cfg'
\quad\Longleftrightarrow\quad
\Canon(\Cfg) \Rightarrow^{*}_{B,D} \Canon(\Cfg').
\]
\end{theorem}
\noindent\textbf{Proof sketch.}
Use Lemma~\ref{lem:alpha-invariance-analysis} to transport each step along an $\alpha$-renaming,
and the fact that $\Canon$ is constant on $\alpha$-classes. \hfill $\square$

\subsection{Finite canonical reachable set}
\label{sec:analysis:finiteness2}

\begin{assumption}[Finite instantiation branching]
\label{ass:finite-inst}
For every $t,\Cfg,\chi$, the set $\SolveGuard(t,\Cfg,\chi)$ is finite.
\end{assumption}

\begin{theorem}[Finiteness of canonical bounded reachability]
\label{thm:finite-canon-reach}
Under Assumption~\ref{ass:finite-inst}, the set
\[
\widehat{V}_{B,D}
~:=~
\left\{\ \Canon(\Cfg)\ \middle|\ \langle M_0,\epsilon\rangle \Rightarrow^{*}_{B,D} \Cfg\ \right\}
\]
is finite.
\end{theorem}
\noindent\textbf{Proof sketch.}
$P$ is finite and budgets bound each $|M(p)|$; stacks are bounded by $D$ over a finite alphabet $\Gamma$.
Canonicalization quotients away unboundedly many names.
Finite branching of $\SolveGuard$ prevents infinite fan-out per node during exploration. \hfill $\square$

\subsection{Monomorphized reachability graph}
\label{sec:analysis:mrg}

We now define the finite reachability graph of canonical bounded semantics.
Edges are labeled by instantiated firings, hence already carry monomorphization information needed for emission.

\begin{definition}[Canonical monomorphized reachability graph]
\label{def:mrg}
Let $\widehat{\Cfg}_0:=\Canon(\langle M_0,\epsilon\rangle)$.
The bounded monomorphized reachability graph (MRG) is the directed labeled graph
$\widehat{G}_{B,D}=(\widehat{V}_{B,D},\widehat{E}_{B,D})$, where $\widehat{V}_{B,D}$ is as in
Theorem~\ref{thm:finite-canon-reach}, and
\[
(\widehat{\Cfg},\varphi,\widehat{\Cfg}')\in \widehat{E}_{B,D}
\quad\Longleftrightarrow\quad
\exists \Cfg,\Cfg'.\
\Canon(\Cfg)=\widehat{\Cfg}\ \wedge\
\Canon(\Cfg')=\widehat{\Cfg}'\ \wedge\
\Cfg[\varphi\rangle_{B,D}\Cfg'.
\]
\end{definition}

\paragraph{Graph construction with parents.}
We build $\widehat{G}_{B,D}$ by worklist saturation and store one parent edge per discovered node.
This supports witness extraction by backtracking.

\begin{algorithm}[t]
\caption{$\mathsf{BuildMRGWithParents}(\mathcal{N}_{\mathcal{C}},B,D)$}
\label{alg:buildmrg-bounded}
\begin{algorithmic}[1]
\REQUIRE $\widehat{V}\gets \emptyset,\ \widehat{E}\gets \emptyset,\ \mathsf{parent}(\cdot)\gets \bot$
\STATE $\widehat{\Cfg}_0\gets \Canon(\langle M_0,\epsilon\rangle)$
\STATE $\mathsf{W}\gets \{\widehat{\Cfg}_0\}$
\WHILE{$\mathsf{W}\neq \emptyset$}
  \STATE remove some $\widehat{\Cfg}=\langle \widehat{M},\widehat{S}\rangle$ from $\mathsf{W}$
  \IF{$\widehat{\Cfg}\in \widehat{V}$} \textbf{continue} \ENDIF
  \STATE $\widehat{V}\gets \widehat{V}\cup\{\widehat{\Cfg}\}$
  \FORALL{$t\in T$}
    \FORALL{choices $\chi$ for ${}^{\bullet}t$ at $\widehat{M}$}
      \FORALL{$\iota\in \SolveGuard(t,\widehat{\Cfg},\chi)$}
        \STATE $\varphi \gets \langle t,\chi,\iota\rangle$
        \STATE $\Cfg'\gets \Fire(\widehat{\Cfg},\varphi)$
        \IF{$\mathsf{BudOK}_{B,D}(\Cfg')$}
          \STATE $\widehat{\Cfg}'\gets \Canon(\Cfg')$
          \STATE $\widehat{E}\gets \widehat{E}\cup\{(\widehat{\Cfg},\varphi,\widehat{\Cfg}')\}$
          \IF{$\widehat{\Cfg}'\notin \widehat{V}\ \AND\ \mathsf{parent}(\widehat{\Cfg}')=\bot$}
            \STATE $\mathsf{parent}(\widehat{\Cfg}')\gets (\widehat{\Cfg},\varphi)$
          \ENDIF
          \STATE $\mathsf{W}\gets \mathsf{W}\cup\{\widehat{\Cfg}'\}$
        \ENDIF
      \ENDFOR
    \ENDFOR
  \ENDFOR
\ENDWHILE
\RETURN $(\widehat{G}_{B,D}=(\widehat{V},\widehat{E}),\mathsf{parent})$
\end{algorithmic}
\end{algorithm}

\begin{theorem}[Termination and completeness within bounds]
\label{thm:mrg-bounded-term-comp}
Under Assumption~\ref{ass:finite-inst}, Algorithm~\ref{alg:buildmrg-bounded} terminates and returns
the complete MRG $\widehat{G}_{B,D}$, and a parent map that yields a witness firing sequence
to every node in $\widehat{V}_{B,D}$.
\end{theorem}
\noindent\textbf{Proof sketch.}
By Theorem~\ref{thm:finite-canon-reach} the canonical reachable set is finite.
For each node, the number of choices $\chi$ is finite (bounded multiplicities), and by Assumption~\ref{ass:finite-inst}
each $\SolveGuard$ call returns a finite set of instantiations. Hence the worklist saturation terminates.
Completeness holds because the algorithm enumerates \emph{all} bounded-enabled instantiated firings from every discovered node.
\hfill $\square$

\subsection{Classical decision problems on $\widehat{G}_{B,D}$}
\label{sec:analysis:classic2}

Since $\widehat{G}_{B,D}$ is finite, standard analysis problems reduce to graph algorithms.

\paragraph{Reachability with witness extraction.}
Let $\Phi$ be an objective predicate over canonical configurations.
We say $\Phi$ is \emph{achievable within bounds} iff
$\exists \widehat{\Cfg}\in \widehat{V}_{B,D}.\ \Phi(\widehat{\Cfg})$.
A witness firing sequence is obtained by backtracking $\mathsf{parent}$ from such a node.

\begin{theorem}[Bounded reachability]
\label{thm:bounded-reachability}
Assume Assumption~\ref{ass:finite-inst}. For any predicate $\Phi$,
$\Phi$ is achievable within bounds iff Algorithm~\ref{alg:buildmrg-bounded} discovers some node
$\widehat{\Cfg}\in\widehat{V}_{B,D}$ with $\Phi(\widehat{\Cfg})$.
\end{theorem}
\noindent\textbf{Proof.}
Soundness: every discovered node is reachable by construction of edges from bounded firings.
Completeness: every bounded-reachable canonical node will eventually be saturated and inserted into $\widehat{V}_{B,D}$.
\hfill $\square$

\subsection{Searching $\widehat{G}_{B,D}$ for compilable witnesses}
\label{sec:analysis:classic2}

% 可达图构建 + 可达图上的搜索策略
% 在有限可达图 \widehat{G}_{B,D} 上“如何搜索到一个我们想要的 firing 序列”，
% 并且明确两类策略痛点：
%   (A) const/值生成器（含 Copy/Clone 的复制规则或零参构造器）可反复发生，会把探索空间刷爆；
%   (B) 变迁触发时的参数绑定/泛型实例化往往多解，需要策略选择“更好”的绑定与类型。
%
% 不论采用何种策略，我们都要求每一步仍必须通过 SolveGuard 才能触发；
% 所以正确性不受影响。策略只影响“探索顺序/剪枝/采样”。

Since $\widehat{G}_{B,D}$ is finite, code-generation objectives reduce to graph search. Let $\Phi$ be an objective predicate over canonical configurations.

% 这段保留“理论完备性”版本：如果你确实想把可达图全量构建出来，Alg~\ref{alg:buildmrg-bounded} 给出 complete graph。
We say $\Phi$ is achievable within bounds iff $\exists \widehat{\Cfg}\in \widehat{V}_{B,D}.\ \Phi(\widehat{\Cfg})$. A witness firing sequence can be obtained by backtracking $\mathsf{parent}$ from such a node.

\begin{theorem}[Bounded reachability via full exploration]
\label{thm:bounded-reachability}
Assume Assumption~\ref{ass:finite-inst}. For any predicate $\Phi$,
$\Phi$ is achievable within bounds iff Algorithm~\ref{alg:buildmrg-bounded} discovers some node
$\widehat{\Cfg}\in\widehat{V}_{B,D}$ with $\Phi(\widehat{\Cfg})$.
\end{theorem}
\noindent\textbf{Proof.}
Soundness: every discovered node is reachable by construction of edges from bounded firings. Completeness: every bounded-reachable canonical node will eventually be saturated and inserted into $\widehat{V}_{B,D}$.
\hfill $\square$

\paragraph{Why strategies matter (goal-directed generation).}
% 这一段对应你提的两个策略关键。
Even though the bounded graph is finite, exhaustive exploration can still be wasteful for generation. Two phenomena are especially common:
(i) endlessly applicable producers: transitions that can repeatedly create additional values (e.g., \texttt{const}-like
constructors, zero-input value seeds, or duplication rules for \texttt{Copy}/\texttt{Clone}) may dominate the frontier;
(ii) binding explosion: a single transition schema $t$ may admit many token choices $\chi$ and many instantiations $\iota$
when multiple ground types satisfy the same obligations.

\paragraph{A strategy interface.}
% 把“搜索策略”显式化成一组可替换的 policy：
%   - Cost: 代价函数（惩罚 const-like/惩罚深栈 push 等）
%   - PickChoices: 从所有 token 绑定 χ 中选子集（可 top-k / 采样）
%   - PickInst: 从 SolveGuard 的多解中选子集（可 top-k）
%   - RankTy: 对 completion 时的候选类型排序（复用已有类型、偏好更简单类型等）
A search strategy $\mathsf{Strat}$ provides:
(a) a nonnegative cost function $\mathsf{Cost}(\varphi)$ for instantiated firings;
(b) a policy $\mathsf{PickChoices}(\widehat{\Cfg},t)$ that selects a finite subset of token choices $\chi$;
(c) a policy $\mathsf{PickInst}(\widehat{\Cfg},t,\chi)$ that selects a finite subset of instantiations from $\SolveGuard$.

A simple and effective cost design is to penalize producer-like steps and deepening borrows:
\[
\mathsf{Cost}(\langle t,\chi,\iota\rangle)
~=~
1
~+~
w_{\mathrm{prod}}\cdot \delta_{\mathrm{prod}}(t)
~+~
w_{\mathrm{push}}\cdot \delta_{\mathrm{push}}(t),
\]
where $\delta_{\mathrm{prod}}(t)\in\{0,1\}$ flags \texttt{const}-like/value-producing transitions (including duplication rules),
$\delta_{\mathrm{push}}(t)\in\{0,1\}$ flags transitions whose action is a stack push, and $w_{\mathrm{prod}},w_{\mathrm{push}}$ are
user-chosen weights.
% 生成目标通常需要有效的 API 调用序列, 过多的 producer 会让序列变长但不推进目标；过深的 borrow 栈会让 closing 更难。

\paragraph{Goal-directed exploration (on-the-fly).}
% 1) 算法并不要求先把整个 \widehat{G}_{B,D} 全量建完；
%    边搜索边展开的 best-first / Dijkstra 。
% 2) 解决const 无限发”问题：通过 Cost 惩罚 + PickChoices/PickInst 限制分支，
%    搜索会优先走更可能推进目标的边。
% 3) 正确性来源于：每条边仍然来自 SolveGuard + bounded firing，因此输出的序列一定可编译（在模型意义下）。

\begin{algorithm}[t]
\caption{$\mathsf{GuidedSearch}(\mathcal{N}_{\mathcal{C}},B,D,\Phi,\mathsf{Strat})$}
\label{alg:guided-search}
\begin{algorithmic}[1]
\REQUIRE $\widehat{\Cfg}_0\gets \Canon(\langle M_0,\epsilon\rangle)$
\STATE priority queue $\mathsf{PQ}\gets \{(0,\widehat{\Cfg}_0)\}$
\STATE $\mathsf{dist}(\widehat{\Cfg}_0)\gets 0$; for all others $\mathsf{dist}(\cdot)\gets +\infty$
\STATE $\mathsf{parent}(\cdot)\gets \bot$
\WHILE{$\mathsf{PQ}\neq \emptyset$}
  \STATE pop $(d,\widehat{\Cfg})$ with minimal $d$ from $\mathsf{PQ}$
  \IF{$d>\mathsf{dist}(\widehat{\Cfg})$} \textbf{continue} \ENDIF
  \IF{$\Phi(\widehat{\Cfg})$} \RETURN $(\widehat{\Cfg},\mathsf{parent})$ \ENDIF
  \FORALL{$t\in T$}
    \FORALL{$\chi\in \mathsf{PickChoices}(\widehat{\Cfg},t)$}
      \FORALL{$\iota\in \mathsf{PickInst}(\widehat{\Cfg},t,\chi)$}
        \STATE $\varphi\gets \langle t,\chi,\iota\rangle$
        \STATE $\Cfg'\gets \Fire(\widehat{\Cfg},\varphi)$
        \IF{$\mathsf{BudOK}_{B,D}(\Cfg')$}
          \STATE $\widehat{\Cfg}'\gets \Canon(\Cfg')$
          \STATE $d'\gets d+\mathsf{Cost}(\varphi)$
          \IF{$d'<\mathsf{dist}(\widehat{\Cfg}')$}
            \STATE $\mathsf{dist}(\widehat{\Cfg}')\gets d'$
            \STATE $\mathsf{parent}(\widehat{\Cfg}')\gets (\widehat{\Cfg},\varphi)$
            \STATE push $(d',\widehat{\Cfg}')$ into $\mathsf{PQ}$
          \ENDIF
        \ENDIF
      \ENDFOR
    \ENDFOR
  \ENDFOR
\ENDWHILE
\RETURN $\bot$
\end{algorithmic}
\end{algorithm}

\paragraph{Soundness of strategy-guided witnesses.}
% 这段结论对应“策略不会破坏可编译性”，只影响探索的充分性/效率。
Any firing sequence obtained by backtracking $\mathsf{parent}$ from a returned goal node consists only of
bounded-enabled firings (each step comes from $\SolveGuard$ and respects $\mathsf{BudOK}_{B,D}$),
hence it is a valid bounded witness in the PCPN semantics. Completeness is recovered by choosing
$\mathsf{PickChoices}$ and $\mathsf{PickInst}$ to enumerate all possibilities (reducing to full exploration).

\subsection{Closed witnesses and self-contained snippet synthesis}
\label{sec:analysis:synth2}

We call a configuration \emph{closed} iff its stack is empty:
\[
\Closed(\langle M,S\rangle)\ \triangleq\ S=\epsilon.
\]
A firing sequence is \emph{closed} iff it reaches a closed configuration.

\paragraph{Deterministic stack discharge.}
We close open scopes by repeatedly discharging the topmost stack frame.
Let $\mathsf{EndTransitions}(\gamma)$ return the finite set of ending schemas compatible with the frame kind:
e.g., for $\mathsf{Mut}(\cdot)$ use $\mathsf{EndMut}$ or $\mathsf{EndProjMut}$ depending on the stored ids/types;
for $\mathsf{Shr}(\cdot)$ use $\mathsf{EndShrKeep}$ or $\mathsf{EndShrLast}$; for $\mathsf{Freeze}(\cdot)$
it is discharged together with the matching last shared-borrow end.

\begin{algorithm}[t]
\caption{$\mathsf{CloseStack}(\Cfg)$}
\label{alg:closestack2}
\begin{algorithmic}[1]
\REQUIRE configuration $\Cfg=\langle M,S\rangle$
\ENSURE closing suffix $\sigma_{\mathrm{close}}$ or $\bot$
\STATE $\sigma_{\mathrm{close}}\gets \epsilon$
\WHILE{$S\neq \epsilon$}
  \STATE $\gamma \gets \mathsf{top}(S)$
  \STATE $E \gets \mathsf{EndTransitions}(\gamma)$
  \STATE $\textbf{found}\gets \textbf{false}$
  \FORALL{$t\in E$}
    \FORALL{choices $\chi$ for ${}^{\bullet}t$ at $M$}
      \FORALL{$\iota\in \SolveGuard(t,\Cfg,\chi)$}
        \STATE $\varphi\gets \langle t,\chi,\iota\rangle$
        \STATE $\Cfg\gets \Fire(\Cfg,\varphi)$
        \STATE $\sigma_{\mathrm{close}}\gets \sigma_{\mathrm{close}}\cdot \varphi$
        \STATE $\textbf{found}\gets \textbf{true}$
        \STATE \textbf{break all loops}
      \ENDFOR
    \ENDFOR
  \ENDFOR
  \IF{$\textbf{found}=\textbf{false}$} \RETURN $\bot$ \ENDIF
\ENDWHILE
\RETURN $\sigma_{\mathrm{close}}$
\end{algorithmic}
\end{algorithm}

Given a witness $\sigma$ reaching $\Cfg_\sigma$, define the closed extension
$\sigma^{\downarrow} := \sigma\cdot \mathsf{CloseStack}(\Cfg_\sigma)$ when $\mathsf{CloseStack}\neq\bot$.

\begin{lemma}[Closing correctness]
\label{lem:closing-correct}
If $\mathsf{CloseStack}(\Cfg)=\sigma_{\mathrm{close}}\neq\bot$, then
$\Cfg[\sigma_{\mathrm{close}}\rangle \Cfg'$ for some $\Cfg'$ with $\Closed(\Cfg')$.
\end{lemma}
\noindent\textbf{Proof sketch.}
Each iteration fires an ending transition that pops at least one top frame; thus the stack strictly decreases.
Guards enforce correct matching and restoration of resources, so the process terminates exactly at empty stack. \hfill $\square$

\paragraph{Emission soundness.}
We keep the emission interface minimal: each instantiated firing emits one BSNF statement
(call/borrow/projection/drop) with exactly the monomorphization recorded in $\iota$.

\begin{assumption}[Stepwise emission correctness]
\label{ass:emit-stepwise}
For every enabled instantiated firing $\varphi$ at configuration $\Cfg$,
the statement $\mathrm{EmitCore}(\varphi)$ is accepted by the Rust type checker and borrow checker
under $\Sigma(\mathcal{C})$, and its resource effect matches $\Fire(\Cfg,\varphi)$.
\end{assumption}

\begin{theorem}[Soundness of emitted closed witnesses]
\label{thm:emit-sound}
Assume Assumption~\ref{ass:emit-stepwise}.
If $\sigma^{\downarrow}$ is a closed firing sequence in the bounded semantics,
then the concatenation $\mathrm{EmitCore}(\sigma^{\downarrow})$ is a self-contained Rust snippet
accepted under $\Sigma(\mathcal{C})$.
\end{theorem}
\noindent\textbf{Proof sketch.}
By Assumption~\ref{ass:emit-stepwise}, each emitted statement is accepted and preserves the modeled resource discipline.
By Lemma~\ref{lem:closing-correct}, the appended suffix discharges all outstanding borrows, so the snippet is closed/self-contained.
Hence the full snippet is accepted. \hfill $\square$

\subsection{Progressive completeness under increasing bounds}
\label{sec:analysis:progressive2}

Global completeness over an unbounded type universe is impossible.
However, every \emph{finite} accepted BSNF program uses finitely many monomorphic types,
allocates finitely many values, and has a finite maximum borrow nesting depth.
Thus it lies within some finite bound instance $(B,D)$ and finite type universe $\widehat{\mathsf{Ty}}$.

\begin{theorem}[Limit completeness for finite witnesses]
\label{thm:limit-completeness}
Let $P$ be any BSNF program accepted by the SRS under $\Sigma(\mathcal{C})$.
Then there exist finite bounds $B,D$ and a finite ground-type universe $\widehat{\mathsf{Ty}}$
such that $P$ corresponds to a PCPN firing sequence in $\Rightarrow^{*}_{B,D}$ reaching a closed configuration.
Consequently, an iterative deepening strategy that increases $(\widehat{\mathsf{Ty}},B,D)$ will eventually discover a witness for $P$.
\end{theorem}
\noindent\textbf{Proof sketch.}
Let $\widehat{\mathsf{Ty}}$ contain all monomorphic types appearing in $P$.
Let $D$ be the maximum nested borrow depth in $P$.
Let $B(p)$ bound the maximum number of simultaneously live variables of place-kind $p$ during $P$'s execution.
Then $P$ fits within $(\widehat{\mathsf{Ty}},B,D)$ by construction. \hfill $\square$

\section{Experiment}
% 有必要加吗哥们？

\section{Related Work}
% ==========================
\subsection{Positioning: type-only Petri nets, LLMs, and our compilation-aligned model}
\label{sec:positioning}
\subsection{Petri Nets}
% SyPet 的关键点：place=type, transition=method, token=该类型变量数量；
% reachability 找到方法序列，但这只是 program sketch（参数绑定未知），
% 还需要 sketch completion（约束+SAT）来决定“每个洞到底用哪个变量/常量”。
SyPet~\cite{FengPOPL17SyPet} models relationships between library methods using a compact Petri-net representation that places correspond to types, transitions represent methods, and tokens abstract the number of available variables of each type. Reachability yields a sequence of method calls; however, the reachable path corresponds to a program sketch rather than a complete program, and a separate completion phase is required to choose actual arguments for each call.
% SyPet抽象层级不同导致缺口对于 Safe Rust
% 1) token 只表示数量（多重集）→变量身份/别名关系被抹平→参数绑定必须在第二阶段做；
% 2) Rust 的关键难点是 capability + borrow/lifetime 的演化
%    这些不是类型到类型的静态流能表达的，需要把借用栈和能力变化进状态；
% 3) SyPet 论文也讨论过一种更精确的网（节点代表参数/返回值）能直接翻译成代码，
%  但它会显著增大网并使 reachability 更难扩展——这正说明只靠类型层 Petri net会在精确性与规模间两难。
For Safe Rust, compilation feasibility depends not only on the compatibility of input and output types, but also on several interrelated semantic constraints: the correct transfer of ownership for non‑Copy values, the enforcement of aliasing discipline under mutable references, the preservation of memory immutability under shared borrows, and the proper scoping of permissions through lifetimes. A Petri net that abstracts program states using only type multiplicities necessarily discards variable identity and the structure of the borrow stack. Consequently, critical decisions—such as argument selection, alias management, and the scheduling of permission discharge—must be deferred to a separate post-synthesis completion phase. SyPet explicitly acknowledges that adopting a more precise Petri-net encoding, one that directly maps execution runs to executable code, would substantially increase the size of the net and render the reachability problem significantly less scalable. Our PCPN bridges this gap by elevating borrow-stack state and capability evolution to first-class elements of the net’s state space. As a result, reachability in the PCPN directly corresponds to compilation-feasible program traces, modulo our assumptions regarding BSNF and encoding fidelity.

\subsection{Large Language Models}
% 1) LLM 缺少可证明的可编译性保证（尤其 Rust borrow checker）；
% 2) 上下文窗口有限：大 crate 的签名、trait bound、关联类型等塞不进去→容易遗忘约束；
% 3) LLM 的搜索不是系统性的：容易陷入局部模式，遇到需要长轨迹/精确资源管理就不稳定；
% 4) LLM 输出具有随机性：同一提示多次生成不一致→在 fuzz/repair 中难复现。
Recent testing and repair workflows frequently leverage large language models to propose candidate code. However, LLM-generated outputs offer no guarantee of compilation feasibility—particularly under Rust’s borrow checker, which enforces intricate ownership and aliasing rules. Compounding this challenge, the effective context window of current models restricts how many function signatures, trait bounds, and associated-type constraints can be presented simultaneously, raising the risk of omitting essential obligations or introducing contradictory ones. Furthermore, LLM-driven code generation relies on heuristic exploration rather than systematic search, often revisiting unproductive patterns—such as repeatedly applying constructor-like calls—without mechanisms for state-aware pruning. For these reasons, we treat LLMs not as primary search engines, but as completers that operate within a rigorously constrained space.
Concretely, our reachability search yields a witness firing sequence $\sigma$ along with a finite type pool $\widehat{\mathsf{Ty}}$ and per-hole typing constraints. The LLM is then tasked with synthesizing surface-level syntax—including qualified paths, minor rewrites, constants, and small expressions—that satisfy these constraints. Compilation feedback can subsequently guide localized revisions to the filled holes, while the underlying PCPN witness preserves the global structure of the feasible execution trace.

\subsection{Others}
% 其它工具要么需要强规格/逻辑规范，要么规模不行，要么对 Rust borrow/lifetime 没对齐语义。
Existing component-based or enumerative synthesis tools either require rich semantic specifications, do not scale to long stateful traces, or do not model Rust-specific compilation constraints in the search state. 


\section{Conclusion}
\label{sec:conclusion}
This paper studied that given an extracted environment $\Sigma(\mathcal{C})$, we asked which long, stateful call traces can be realized by a compilable client without inspecting function bodies. Our approach is semantics-first and treats compiler-enforced signature-level constraints as the notion of trace realizability.

We made three main contributions. First, we formulated a SRS whose configurations are multisets of capability-indexed typed resources together with a loan stack, and whose steps are induced solely by signature information and a fixed family of structural rules. Second, we constructed a crate-indexed PCPN directly from $\Sigma(\mathcal{C})$, including judgemental unification and obligation checking at firing time, and proved a strong bisimulation between the SRS LTS and the PCPN firing LTS (Theorem~\ref{thm:bisim}). This yields an if-and-only-if characterization of realizable monomorphized traces as those labeling firing sequences that reach a closed configuration. Third, under explicit bounds on token multiplicities and borrow depth, we defined $\alpha$-canonicalization and
obtained a finite quotient reachability graph that serves as a search space for witness extraction and snippet emission; we further described strategy-guided exploration to throttle endlessly applicable value producers and to choose effective parameter/type bindings when multiple instantiations match.

Our results are deliberately conservative and scoped. We focus on the safe fragment and signature-level constraints; we do not model unsafe code, interior mutability patterns, or semantic protocols that depend on function bodies.
Moreover, our pushdown encoding assumes well-bracketed borrow scopes in the emitted core fragment, and completeness claims are relative to explicit finiteness budgets and finite instantiation branching.

Future work includes relaxing the well-bracketed restriction to approximate more of Rust's NLL behavior while retaining analyzability, strengthening and mechanizing the link between SRS admissibility and \texttt{rustc} acceptance for the chosen core fragment, improving generic instantiation strategies and trait/associated-type solving in large extracted environments, and scaling an implementation to real-world crates and integrating witness emission with
debugging-oriented counterexample explanations.


\begin{figure}
    \centering
    \includegraphics[width=\linewidth]{graphviz.png}
    \caption{Caption}
    \label{fig:placeholder}
\end{figure}
\end{document}
